\section{拓扑空间}

\noindent\textbf{1. 拓扑空间的基本概念}\quad
度量空间是最直观的拓扑空间，它通过两点间的距离来定义开集、闭集、邻域、收敛性等等，与微积分中的讲法很接近。而且在这个基础上我们特别建立了巴拿赫空间和希尔伯特空间。那么为什么还要进一步抽象化呢？一方面固然是因为确有不可定义度量的空间（然而这类空间并不常见，有许多是属于“病态空间”之列）；另一方面，有许多情况下使用度量并不方便。这样就使得我们进一步考虑，在度量“后面”是否还有更深一层的结构？我们弄清楚了这些结构，就能更深刻地理解微积分的一些概念、方法之实质。例如在 \S1 中关于紧性的问题就有一个未证明的定理，讲到紧和列紧——即一个序列是否有收敛子序列的问题，又如在那里讲到连续函数的一些性质，我们只说关于连续函数可以取中间值的定理和其它两个性质（达到最大最小值以及一致连续性）是不同的，它涉及连通性。那么什么是连通性呢？我们学过的微积分知识中还有哪一些是涉及连通性的呢？因此，我们的目的并不是要去研究一些很特殊的空间。虽然我们主张，如果一个空间很有用，就应该去讨论它，哪怕不太容易也要讨论。我们的目的在于探究一些更深刻的结构，以便用来解决更多的问题。

举一个例子，设函数 $y=f(x)$ 在 $x=x_0$ 点连续，其定义是任给一个 $\varepsilon>0$，必有相应的 $\delta(\varepsilon)>0$ 存在，使当 $|x-x_0|<\delta$ 时，$|y-f(x_0)|<\varepsilon$。如果用邻域的语言来陈述就是：设有 $f(x_0)$ 的一个邻域 $V$，必可找到 $x_0$ 的相应邻域 $U$，使得 $f(U)\subset V$。这里，$\varepsilon$ 和 $\delta$ 分别是用以度量 $V,U$ 的两个数。也就是说，邻近的程度是用数来刻画的。现在要问，如果不用数的大小可否仍然讨论“邻近”的概念？当然这样做会产生一些困难，例如 $f(x_0)$ 的邻域 $V$ 若用 $|y-f(x_0)|<\varepsilon$ 来刻画，则两个邻域可以“比较大小”，可以讨论一串邻域逐渐趋向一点等等，如果没有度量，这些概念如何刻画？把这一类问题暂时放在一边，应该说，“邻近”的概念是否一定要用度量来刻画是一个问题。

如果不用度量来刻画，就只有用公理化的方法来“定义”什么样的集合可以说是 $x_0$ 或 $f(x_0)$ 的邻域，例如至少应该要求 $x_0\in(x_0\text{ 的任一邻域})\cdots\cdots$。

认识到有必要这样来讨论空间中有关连续性的问题或称为拓扑问题，经历了一个漫长的时期，可以追溯到19世纪后半叶，其奠基者应该算是黎曼。到了20世纪初，人们才得到了共识，有了拓扑空间的理论。这个理论采取了公理化的形式，也是那个时代数学的总的潮流决定的。对于一个空间的拓扑性质，即有关连续性的性质，可以从一个或几个基本概念开始讲起，至于选取哪一个概念，有不同的作法。有主张从邻域开始，有主张从闭包开始，而后来大家公认的讲法则是从开集开始。这段历史我们不能细说了，只是要提醒，当年拓扑学的发展是为了更深刻地研究许多数学和物理学问题——一个重大的问题是天体力学问题，而不只是为了对数学概念作更精确的描述。

下面我们就按拓扑学发展的这种精神，从最基本的概念开始作公理化的叙述。

\noindent\textbf{定义1}\quad
空间（即一个集合）$X$ 之一个拓扑 $\tau$ 就是指 $X$ 的一族子集 $\mathcal O$，其中之元（都是 $X$ 的子集）称为开集，$\tau$ 应该适合以下的条件：
\begin{enumerate}[label=(\roman*)]
\item $X$ 和 $\varnothing$ 都是开集：$X\in\tau,\ \varnothing\in\tau$；
\item 任意多个 $\tau$ 中之元，即任意多个开集之并仍为开集：$\bigcup_iU_i\in\tau\ (U_i\in\tau)$；
\item 有限多个 $\tau$ 中之元，即有限多个开集 $U_1,\ldots,U_N$ 之交仍为开集，即 $\bigcap_{i=1}^NU_i\in\tau$。
\end{enumerate}
$X$ 赋以这样的拓扑后称为一个拓扑空间。

上一节中我们看到在度量空间 $X$ 中，以球为基础，定义了开集。我们还特别约定空集是开集，全空间是开集则很容易理解。球是一个与度量（半径长）有关的概念，所以那里的开集并非是先验地加以规定的，它是否有现在的定义中的(i)--(iii)的性质，是应该证明的。我们也确实证明了它们（\S2 定理2）。这样我们看到，在度量空间中按 \S2 定义4所定义的开集，确实符合现在所给的开集的定义。于是读者会问，按现在的定义所给的开集，对于度量空间，是否恰好就是 \S2 定义4 的开集呢？这是不一定的。因为 \S2 定义4 如上所述是以度量为基础的，而这里的定义完全不涉及度量。从完全不涉及度量的定义出发，要能推导出基于度量的概念，是很不容易的。一定要对 $X$ 加上一些限制才行。这时，我们就说 $X$ 可度量化（metrizable）。而且重要的是，同一个集合 $X$（我们暂时不用空间这个词）可以赋予完全不同的拓扑，构成完全不同的空间，怎么可能都是度量空间呢？大家都会看到，$X$ 上至少有以下两种拓扑：

\begin{enumerate}[label=\roman*)]
\item $(X,\tau_0)$（$\tau_0$ 就是指适合定义1的(i)--(iii)的开集族或者称为一个拓扑），这里 $\tau_0=\{\varnothing,X\}$，它显然适合定义1之要求。$\tau_0$ 称为平凡拓扑。
\item $(X,\tau_1)$，$\tau_1=\{X\text{ 之所有子集}\}$。当然 $\varnothing$ 和 $X$ 都是 $X$ 的子集。$\tau_1$ 称为离散拓扑。
\item 如果可能在 $X$ 上引入度量 $\rho$，则可以用这个度量来定义 \S2 定义4 那样的拓扑，因此，这种拓扑称为度量拓扑。
\end{enumerate}

于是我们看到，任给一个集合 $X$，都至少可以赋以 $\tau_0$ 与 $\tau_1$ 两种拓扑。但判断一个集合是否可以赋以度量拓扑则是很困难的事。这样我们也就看到，度量空间只是一种很特殊的拓扑空间。因为同一个集合 $X$ 可以赋予不同的拓扑而成不同的拓扑空间，所以有时我们也把有关的拓扑 $\tau$ 写出来而记一个拓扑空间为 $(X,\tau)$。

同一集合 $X$ 可以赋予不同的拓扑，这就产生了拓扑的比较问题。设同一集合 $X$ 上有两个拓扑 $\tau$ 与 $\tau'$，即指定两个子集族满足定义1的各项要求，并分别称其中之元为 $\tau$ 开集与 $\tau'$ 开集。若一切 $\tau$ 开集均为 $\tau'$ 开集，即由 $U\in\tau$ 可得 $U\in\tau'$，就说 $\tau$ 弱于（粗于）$\tau'$，同时也就说 $\tau'$ 强于（细于）$\tau$。所以，平凡拓扑弱于一切拓扑，而离散拓扑则强于一切拓扑。但是并非任意两个拓扑都可以比较强弱。

\noindent\textbf{定义2}\quad
设 $(X,\tau)$ 为一拓扑空间，$U\subset X$ 称为这个拓扑空间（或简单就说是“这个拓扑”）的闭集，如果其余集 $U^c=X\setminus U$ 是一个开集。

闭集有下面的性质（证明很容易，故略去）：

\noindent\textbf{定理1}\quad
若 $(X,\tau)$ 之闭集族记为 $\mathcal F$，则
\begin{align*}
 &(i)\quad X,\varnothing\in\mathcal F;\\
 &(ii)\quad \text{有限多个闭集之并仍为闭集，即若 }U_1,\ldots,U_N\in\mathcal F,
 \text{ 则 }\bigcup_{i=1}^N U_i\in\mathcal F;\\
 &(iii)\quad \text{任意多个闭集之交仍为闭集，即若 }U_i\in\mathcal F,
 \text{ 则 }\bigcap_i U_i\in\mathcal F.
\end{align*}

\noindent\textbf{定义3}\quad
设 $(X,\tau)$ 是一个拓扑空间，$x\in X$ 的邻域是指适合以下条件的 $X$ 之子集 $\Omega(x)$：$x\in\Omega(x)$ 而且存在一个开集 $U$：$x\in U\subset\Omega(x)$。

这个定义与 \S2 定义5（度量空间中邻域的定义）是一致的。$x$ 之邻域的集合称为 $x$ 之邻域系，记作 $\mathcal N_x$。

\noindent\textbf{定理2}\quad
设 $(X,\tau)$ 是一个拓扑空间，$\mathcal N_x$ 为 $x\in X$ 之邻域系，则
\begin{align*}
 &(i)\quad \text{若 }U_i\in\mathcal N_x,\text{ 则 }\bigcup_i U_i\in\mathcal N_x;\\
 &(ii)\quad \text{若 }U_1,\ldots,U_N\in\mathcal N_x,\text{ 则 }\bigcap_{k=1}^{N}U_k\in\mathcal N_x;\\
 &(iii)\quad \text{若 }U\in\mathcal N_x,\text{ 则 }x\text{ 必有一个邻域 }\widetilde U,\text{ 使 }U\text{ 为 }\widetilde U\text{ 中任一点的邻域};\\
 &(iv)\quad \text{若 }U\in\mathcal N_x,\ V\supset U,\text{ 则 }V\in\mathcal N_x.
\end{align*}

定理1,2之证明均略去，因为它们与度量空间中的相应定理（\S2 定理3,定理4）完全一样。开邻域的概念，以及关于 $X$ 和 $\varnothing$ 既为开集又为闭集的讨论在这里也适用。

关于开集与邻域的关系，很容易证明

\noindent\textbf{定理3}\quad
设 $(X,\tau)$ 是一个拓扑空间，$A\subset X$ 非空，$A$ 为开集之充分必要条件是：$A$ 是其每一点的邻域。

\noindent\textbf{证}\quad
\textbf{必要性}\quad 是明显的，因为定义3中的 $\Omega(x)$ 与 $U$ 可以都选为 $A$。

\textbf{充分性}\quad 假对任一点 $x\in A$，$A$ 均为其邻域，由定义3必有开集 $U_x$，使 $x\in U_x\subset A$，取 $A$ 中所有的点 $x$，并取 $U_x$ 之并，于是 $\{x\}\subset\bigcup_x U_x\subset A$，亦即 $A\subset\bigcup_x U_x\subset A$，所以 $A=\bigcup_x U_x$。由开集的性质(ii)，$A$ 是开集。

下面我们把度量空间中集合 $A$ 的内域、闭包等概念也都移到拓扑空间 $(X,\tau)$ 中来。这里的结果与 \S2 中的定义6,7和定理5也都是一样的，所以证明全部略去。

\noindent\textbf{定义4}\quad
设 $(X,\tau)$ 为一拓扑空间，$A\subset X,x\in X$。
\begin{enumerate}[label=(\roman*)]
\item $x$ 称为 $A$ 之内点，如果 $x$ 有一个开邻域 $U\subset A$；
\item $x$ 称为 $A$ 之外点，如果 $x$ 是 $A$ 之余集 $A^c$ 之内点；
\item 若 $x$ 既非 $A$ 之内点，又非 $A$ 之外点，则称 $x$ 是 $A$ 的边界点。
\end{enumerate}

\noindent\textbf{定义5}\quad
同上，$A$ 之内点的集合称为 $A$ 之内域，记作 $\mathring A$ 或 $\operatorname{int}A$；$A^c$ 之内域的余集 $[\operatorname{int}(A^c)]^c$ 称为 $A$ 之闭包，记作 $\overline A$；边界点之集合记为 $\operatorname{bd}A=X-[\mathring A\cup\operatorname{int}(A^c)]$，称为 $A$ 的边界。

\noindent\textbf{定理4}\quad
$\mathring A$ 是一切含于 $A$ 中的开集之并，所以是含于 $A$ 内的最大开集；$\overline A$ 是一切包含 $A$ 的闭集之交，所以是包含 $A$ 的最小闭集；$\operatorname{bd}A$ 也是闭集，而且
\begin{equation}
 \overline A=\mathring A\cup\operatorname{bd}A=A\cup\operatorname{bd}A,
 \qquad
 \mathring A=\overline A\setminus\operatorname{bd}A=A\setminus\operatorname{bd}A.
 \tag{1}
\end{equation}
从这个定理还可以看到，$A$ 为开（闭）集的充分必要条件是 $A=\mathring A$（$A=\overline A$）。

我们还要引入一个很有用的概念——导集。为此先给出以下的定义：

\noindent\textbf{定义6}\quad
设 $(X,\tau)$ 为一拓扑空间，$A\subset X,x\in X$ 称为 $A$ 之极限点，如果 $x$ 之任一邻域 $\Omega(x)$ 中均包含有 $A$ 中除 $x$ 以外的其它点，极限点之集合称为 $A$ 之导集，记作 $A'$。

极限点的概念很明显是通常微积分教材中讲到的点的序列之极限点的概念的推广。我们不妨看一下，$x\notin A'$ 是什么意思。由定义6，$x\notin A'$ 显然是指 $x$ 有某一邻域 $\Omega(x)$ 或者与 $A$ 不相交，即其中不含 $A$ 之点，这时 $x$ 是 $A$ 的外点；或者 $\Omega(x)$ 中只有一点 $x\in A$，这时 $x$ 是 $A$ 的孤立点，所以极限点就是在 $X$ 中除去 $A$ 之外点以及 $A$ 之孤立点后余下的点。极限点也称为聚点。

我们在 \S1,\S2 中都讲到了处处稠密的概念，而且看见了它在微积分以至整个分析数学中的重要性。现在我们可以给它一个正式的定义。

\noindent\textbf{定义7}\quad
设 $(X,\tau)$ 是一拓扑空间，$X$ 之子集 $A$ 称为在 $X$ 中稠密，就是指 $\overline A=X$。

在微积分中我们经常会遇到这样一类数学问题：即要在一个拓扑空间的某一子集 $A$ 上讨论一个数学问题。最简单的例子是讨论函数 $f(x_1,\ldots,x_n)$ 在某点 $x_0=(x_1^0,\ldots,x_n^0)$ 的连续性，如果这点是 $f$ 之定义域 $A$ 的内点，自然没有任何困难。如果 $x_0\in\operatorname{bd}A$，则条件 $\lim_{x\to x_0}f(x)=f(x_0)$ 应该改成
\[
 \lim_{\substack{x\to x_0\\x\in A}}f(x)=f(x_0).
\]
如果 $f(x)$ 是一元函数，而 $A$ 是一个区间 $[a,b]$，我们通常都说在其端点应该用左、右连续来代替连续。但若 $f(x_1,\ldots,x_n)$ 是多元函数，而 $A$ 是 $\mathbb R^n$ 的某一子集，则 $\operatorname{bd}A$ 的构造可以非常复杂，这时把 $A$ 本身就当作一个拓扑空间就方便多了。因此，我们宁可用 $(A,\tau')$ 代替 $(X,\tau)$。问题是 $\tau'$ 如何选择，我们当然希望 $\tau'$ 与 $\tau$ 有关。所以我们一定采用所谓子空间拓扑：

\noindent\textbf{定义8}\quad
设 $(X,\tau)$ 为一拓扑空间，而 $A\subset X$，在 $A$ 中定义拓扑如下：$U\subset A$ 为开集，当且仅当 $U=V\cap A$，而 $V$ 是 $(X,\tau)$ 中的开集。$A$ 上的这个拓扑称为子空间拓扑。若记它为 $\tau'$，我们用 $(A,\tau')\subset(X,\tau)$ 表示不仅 $A$ 是 $X$ 的子集合，而且 $\tau'$ 是子空间拓扑（又称为 $X$ 在 $A$ 上诱导的拓扑、诱导拓扑或相对拓扑）。

以下凡说到某拓扑空间是 $X$ 的子空间时，总是指其上赋有子空间拓扑。当然这里应该验证，$\tau'$ 确实是拓扑，即它适合定义1中的(i)到(iii)，这是很容易的，所以留给读者自行验证。

这样一来，关于左、右连续，读者自然会看出，即是在子空间拓扑下的连续性。而且读者不妨证明一下一个简单的事实：若 $x_0$ 是 $A$ 的孤立点，则任一定义在 $A$ 上的函数 $f(x)$ 在 $x_0$ 处都连续。这个例子会帮助读者不必费很大的劲去专门讨论 $\operatorname{bd}A$ 的构造。

拓扑学的主题不但是拓扑空间，还有拓扑空间之间的映射。其中最重要的是两个拓扑空间 $(X,\tau),(Y,\sigma)$ 之间的连续映射。

\noindent\textbf{定义9}\quad
$f:(X,\tau)\to(Y,\sigma)$ 称为连续的，如果 $(Y,\sigma)$ 之每个开集 $V$ 在 $f$ 下的原像 $f^{-1}(V)=U$ 是 $(X,\tau)$ 的开集。

但是我们在微积分教材中习惯的是一个映射（一个函数）在 $x_0$ 点的连续性，而 $f$ 在整个区域上为连续即指 $f$ 在每点上均为连续。在一般拓扑空间中我们也可以这样做，而有

\noindent\textbf{定理5}\quad
$f:(X,\tau)\to(Y,\sigma)$ 为连续，当且仅当 $f$ 在 $X$ 之每一点 $x_0$ 均为连续，即对 $f(x_0)$ 在 $(Y,\sigma)$ 中之每个邻域 $V_{f(x_0)}$ 均可找到 $x_0$ 在 $(X,\tau)$ 中的一个邻域 $U_{x_0}$，使 $f(U_{x_0})\subset V_{f(x_0)}$。

\noindent\textbf{证}\quad
\textbf{(i)必要性}\quad 设 $f$ 为连续，任取一点 $x_0$，并记 $f(x_0)=y_0$，任取 $y_0$ 在 $(Y,\sigma)$ 中的邻域 $V$。自然有 $f^{-1}(V)=U\subset X$，今证 $U$ 必为 $x_0$ 的邻域。首先，$x_0\in U$ 是明显的；其次 $V$ 既是 $y_0$ 之邻域，必有含 $y_0$ 的开集 $V_1\subset V$。显然 $x_0\in f^{-1}(V_1)\subset U$，但由假设，$f$ 是连续的，开集 $V_1$ 之原像 $f^{-1}(V_1)=U_1$ 是 $(X,\tau)$ 中的开集，而 $x_0\in U_1\subset U$，所以由邻域之定义，$U$ 是 $x_0$ 之邻域，而且 $f(U)\subset V$，故 $f$ 在每一点 $x_0$ 均为连续。

\textbf{(ii)充分性}\quad 设 $f$ 在 $X$ 之每一点处均为连续，任取 $(Y,\sigma)$ 中的开集 $V$，今证 $f^{-1}(V)$ 是 $X$ 之开集。如果 $f^{-1}(V)=\varnothing$，$\varnothing$ 当然是 $(X,\tau)$ 中的开集。如果 $f^{-1}(V)\ne\varnothing$，则令 $x_0$ 为 $U=f^{-1}(V)$ 中任一点，$f(x_0)=y_0$ 即成为 $V$ 中之任意点。由邻域之性质，开集 $V$ 是其任意点之邻域，故由假设有 $x_0$ 之邻域 $U_{x_0}$，使 $f(U_{x_0})\subset V$。于是 $U_{x_0}\subset U$，但 $U_{x_0}$ 中应有一个开集包含 $x_0$ 点，我们就用 $U_{x_0}$ 记这个开集，所以可以设 $U_{x_0}$ 为开集。如果 $x_0$ 遍取 $U$ 中各点，则
\[
 U=\bigcup_{x_0\in U}\{x_0\}\subset\bigcup_{x_0\in U}U_{x_0}\subset U.
\]
所以 $U=\bigcup_{x_0\in U}U_{x_0}$，右方是开集 $U_{x_0}$ 之并，所以也是开集。即是说 $f^{-1}(V)=U$ 作为 $V$ 之原像是开集。因此 $f$ 连续。证毕。

读者应该注意，连续性的定义为“开集的原像仍是开集”，而不是“开集的像是开集”。即是说，若 $U$ 是 $(X,\tau)$ 之开集，$f(U)$ 不一定是 $(Y,\sigma)$ 之开集。如果一个映射 $f:(X,\tau)\to(Y,\sigma)$ 映 $X$ 中之一切开（闭）集为开（闭）集，即上述 $f(U)$ 为开（闭）集，这种映射称为开（闭）映射。开（闭）映射的概念在许多数学问题中很重要，但无论如何，它不必是连续映射。我们不来举例了，只是再提醒一下，必须把开映射和连续映射区别开。

在许多数学分支中还会遇到以下的问题：已知 $f:X\to Y$，如果在 $X$ 或 $Y$ 中已赋有了拓扑，问在 $Y$ 或 $X$ 应该给怎样的拓扑才能使 $f$ 为连续？例如已给了拓扑空间 $(X,\tau)$，我们把 $Y$ 中那些原像为 $\tau$ 中开集的子集规定为开集，即令\jiaokan{此处底本以 $V=f(U)$（$U\in\tau$）来定义 $Y$ 中的开集，并使用 $U=f^{-1}(V)$ 及 $f(\bigcap U_k)=\bigcap f(U_k)$。这些等式对一般映射并不成立（前者需额外的单射/饱和条件，后者一般只能得到包含关系）。为保持本段关于“使 $f$ 连续的最强拓扑”的结论成立，改用标准的定义 $\sigma=\{V\subset Y:f^{-1}(V)\in\tau\}$；以下证明相应按原意改正。}
\[
 \sigma=\{V\subset Y:f^{-1}(V)\in\tau\}.
\]
这样 $f:(X,\tau)\to(Y,\sigma)$ 显然是连续的。问题是这样的 $\sigma$ 是否为一拓扑。显然，$\varnothing$ 与 $Y$ 都属于 $\sigma$，因为其原像为 $\varnothing$ 与 $X$，都是 $\tau$ 中开集。若 $V_\lambda\in\sigma$ 对一切 $\lambda$ 成立，而且 $U_\lambda=f^{-1}(V_\lambda)\in\tau$，则
\[
 f^{-1}\!\left(\bigcup_\lambda V_\lambda\right)=\bigcup_\lambda U_\lambda\in\tau,
\]
所以 $\bigcup_\lambda V_\lambda\in\sigma$。此外，对有限多个 $V_k\in\sigma,k=1,2,\ldots,N$，令其原像为 $U_k=f^{-1}(V_k)$，则 $U_k\in\tau$，而
\[
 f^{-1}\!\left(\bigcap_{k=1}^N V_k\right)=\bigcap_{k=1}^N U_k\in\tau,
\]
所以 $\bigcap_{k=1}^N V_k\in\sigma$。因此这样作出的 $\sigma$ 确实是一个拓扑。如果在这样的 $\sigma$ 之外再添加哪怕一个子集 $V$ 作为开集，则按 $V\notin\sigma$ 的定义，其原像 $f^{-1}(V)$ 必不在 $\tau$ 中，所以会破坏 $f$ 的连续性。因此 $\sigma$ 是使 $f$ 连续的最强（细）的拓扑。

与此相似，若在 $Y$ 中已有拓扑 $\sigma$，令 $V\in\sigma$，我们定义 $V$ 在 $f$ 下的原像 $U=f^{-1}(V)$ 为 $X$ 中的开集，和上面一样，连续性定义的开集族 $\tau$ 确实构成一个拓扑 $\tau$，而且 $f:(X,\tau)\to(Y,\sigma)$ 是连续的。但若把 $\tau$ 中的开集减少一个 $U=f^{-1}(V)$，则 $\sigma$ 开集 $V$ 之原像 $U$ 不再为开，因此 $f$ 不再连续。这就是说，要保证 $f$ 为连续，$\tau$ 中的开集一个也不能少，所以定义 $\sigma$ 开集之原像为 $\tau$ 开集，这样得出的拓扑（上面我们已说了这样的 $\tau$ 确可证明为一拓扑）是使 $f$ 为连续最弱（粗）的拓扑。总之，若 $X$ 中已有拓扑 $\tau$，以一切满足 $f^{-1}(V)\in\tau$ 的 $V\subset Y$ 作为 $Y$ 中的开集，可以得到使 $f$ 为连续的最强的拓扑；若 $Y$ 中已有拓扑 $\sigma$，以 $f^{-1}(V),V\in\sigma$ 作为 $X$ 中的开集，可以得到使 $f$ 为连续的最弱的拓扑。

现在得到了一个最重要的概念。

\noindent\textbf{定义10}\quad
设 $f:(X,\tau)\to(Y,\sigma)$ 是(i)一个一一映射，即一对一地把 $X$ 映到整个 $Y$ 上；(ii)$f$ 是连续的；(iii)$f^{-1}$ 也是连续的，则称 $(X,\tau)$ 与 $(Y,\sigma)$ 同胚（homeomorph），$f$ 为一同胚映射（homeomorphism）。

我们用了“一一映射”（bijection）一词，常用的话有单射（injection）与满射（surjection）两个词。前者指 $Y$ 中任一元 $y$ 如果有原像，则原像为唯一的；后者指 $Y$ 中一切点都是 $X$ 中某点的像，就是 $f$ 把 $X$ 映到 $Y$ 上，而不仅只是 $Y$ 内。前者是说方程 $f(x)=y$ 如果有解，则解总是唯一的；后者则说，方程 $f(x)=y$ 对 $Y$ 中的一切 $y$ 都有解（但不一定是唯一的）。一个映射如果既是单射又是满射就称为是一一映射。这时 $f(x)=y$ 对任意 $y\in Y$ 都有唯一的解，所以有逆映射 $f^{-1}:Y\to X$。

同胚的拓扑空间可以看成相同的空间，因为拓扑空间的基本定义就是开集，$X$ 与 $Y$ 同胚表明，不但 $X$ 与 $Y$ 作为点的集合，是通过 $f$ 实现一一对应的，而且 $\tau$ 与 $\sigma$ 作为开集族，也通过同一个 $f$ 实现一一对应，所以粗略地说，$X$ 与 $Y$ 的一切拓扑性质都是以开集为基础来定义的，则在经过同胚映射后，开集是一一对应的，因而 $X$ 中开集的性质也都移到 $Y$ 中的开集上而没有改变。把这一点说明白一点，所谓拓扑性质就是经同胚映射不变的性质。因此，同胚的拓扑空间可以看成是相同的拓扑空间。这一点经过下面的例子就看得很清楚，下面的几个例子在整个数学中起重要的作用，虽然限于2维情况，其重要性并不稍减。

\noindent\textbf{2. 拓扑空间的几个例子}\quad
考虑 $n$ 维单位球面 $S^n$：
\[
 x_1^2+\cdots+x_{n+1}^2=1.
\]
我们只考虑2维球面 $S^2$，并记上半球面（不带边的）为
\[
 S_+^2=\{x\in\mathbb R^3:x_1^2+x_2^2+x_3^2=1,\ x_3>0\}
\]
（图6-3-1），它的闭包 $\overline{S_+^2}$，则是 $S^2\cap\{x_3\ge0\}$，$S_+^2$ 的边界是平面 $x_3=0$ 上的单位圆周，即 $S^1$。我们又记 $D^2$ 为平面 $x_3=0$ 上的单位圆盘：$x_1^2+x_2^2<1$，其边界即 $S^1$，而闭包 $\overline{D^2}$ 是闭圆盘 $x_1^2+x_2^2\le1$。同样可以考虑下半球面 $S_-^2$。在球面 $S^2$ 和平面 $\mathbb R^2:x_3=0$ 上都按我们直觉习惯的方式规定哪一些集合为开集，于是我们得到两个拓扑空间 $S^2$ 与 $\mathbb R^2$。我们要看一下怎样稍作修改即可建立一种同胚关系。

建立这样的同胚有多种方法，有一个最常见的称为球极射影（stereographic projection），其方法是由球面 $S^2$ 的北极 $N(0,0,1)$ 与球面 $S^2$ 上一点 $A$ 联直线，它交 $\mathbb R^2$ 于 $A'$ 点。这里我们取 $\mathbb R^2$ 为赤道平面 $x_3=0$。于是我们让 $A\leftrightarrow A'$ 点。直觉告诉我们，这个对应是一对一的，而且是双方（即指映射 $A\to A'$ 与逆映射 $A'\to A$）都是连续的。现在我们来计算一下。以下凡是上加一撇“$'$”的都代表 $\mathbb R^3$ 中的平面 $\mathbb R^2:x_3=0$ 上之对象，否则代表球面 $S^2:x_1^2+x_2^2+x_3^2=1$ 上之对象，相同的字母代表相对应的对象。我们把 $NS$ 与 $NA$ 所决定的曲面与平面画在图6-3-1中，并设 $O$ 即原点，$ON$ 为 $x_3$ 轴，$OA'$ 为 $x_j$ 轴。这里 $j=1,2$，而我们得到的结果也对 $j=1,2$ 均成立。于是我们有 $ON=1,OA'=x_j',BN=1-x_3,BA=x_j$，由于 $\triangle NOA'\sim\triangle NBA$，所以
\begin{equation}
 x_j'=\frac{x_j}{1-x_3},\qquad j=1,2.
 \tag{2}
\end{equation}

\begin{figure}[htbp]
\centering
\begin{tikzpicture}[scale=0.88, bella figure]
  % left: sphere and equatorial plane
  \draw[thick] (0,0) ellipse (2.15 and 2.15);
  \draw[dashed] (-2.15,0) arc[start angle=180,end angle=360,x radius=2.15,y radius=.55];
  \draw[thick] (2.15,0) arc[start angle=0,end angle=180,x radius=2.15,y radius=.55];
  \draw[thick] (-3.0,-.2)--(3.4,-.2)--(2.2,-1.15)--(-4.0,-1.15)--cycle;
  \coordinate (N) at (0,2.15); \coordinate (S) at (0,-2.15); \coordinate (A) at (1.25,.9); \coordinate (Ap) at (2.8,-.2);
  \draw[thick] (N)--(Ap); \fill (N) circle (1.2pt) node[above] {$N$}; \fill (S) circle (1.2pt) node[below] {$S$};
  \fill (A) circle (1.2pt) node[above right] {$A$}; \fill (Ap) circle (1.2pt) node[right] {$A'$}; \fill (0,0) circle (1.2pt) node[below left] {$O$};
  \node[right] at (3.2,-.75) {$x_3=0$};
  % right: meridian cross section
  \begin{scope}[xshift=7.0cm,yshift=.25cm]
    \draw[thick] (0,-2)--(0,2);
    \draw[thick] (0,2) arc[start angle=90,end angle=-90,x radius=1.65,y radius=2];
    \draw[thick] (-.15,0)--(2.3,0);
    \coordinate (Nr) at (0,2); \coordinate (Sr) at (0,-2); \coordinate (Ar) at (1.25,.55); \coordinate (Apr) at (1.65,0); \coordinate (Br) at (0,.55);
    \draw[thick] (Nr)--(Apr); \draw[thick] (Br)--(Ar);
    \fill (Nr) circle (1.2pt) node[left] {$N$}; \fill (Sr) circle (1.2pt) node[left] {$S$}; \fill (Ar) circle (1.2pt) node[right] {$A$}; \fill (Apr) circle (1.2pt) node[right] {$A'$}; \fill (Br) circle (1.2pt) node[left] {$B$}; \fill (0,0) circle (1.2pt) node[left] {$O$};
  \end{scope}
\end{tikzpicture}
\caption{图6-3-1}
\end{figure}

现求它的逆映射，将(2)平方后对 $j$ 求和有
\[
 \|x'\|^2=\sum_{j=1}^2x_j'^2
 =\sum_{j=1}^2\frac{x_j^2}{(1-x_3)^2}
 =\frac{1-x_3^2}{(1-x_3)^2}
 =\frac{1+x_3}{1-x_3}.
\]
这里我们利用了 $A\in S^2$，所以 $\sum_{j=1}^2x_j^2+x_3^2=1$。由上式得
\begin{equation}
 1-x_3=\frac{2}{1+\|x'\|^2}.
 \tag{3}
\end{equation}
由此即知 $A\leftrightarrow A'$ 把 $\mathbb R^2$ 与除去 $N$ 点（即 $x_3=1$ 处）的球面 $S^2$ 双方一对一且双方连续地对应起来。

注意 $x_3\to1$ 的情况。这时 $\|x'\|\to+\infty$。如果我们认为 $\mathbb R^2$ 上当 $\|x'\|\to+\infty$ 时的点都是一个点，而且认为是无穷远点，或者说把这个“极限”看成一个点添加到 $\mathbb R^2$ 上，就得到“扩充的”$\mathbb R^2$：$\mathbb R^2\cup\{\infty\}$。而且我们就说扩充的 $\mathbb R^2$ 上只有一个无穷远点。于是“扩充的”$\mathbb R^2$ 与 $S^2$ 是互相同胚的，球极射影(2)与(3)就是一个同胚映射。于是按上面所讲的，$\mathbb R^2\cup\{\infty\}$ 与 $S^2$ 应看成相同的拓扑空间。例如 $S^2$ 上北极 $N$ 的邻域就变成扩充的 $\mathbb R^2$ 上无穷远点的邻域。例如，$\|x'\|>R>0$（$R$ 是任意正数）就是无穷远点的一个邻域。

这个观点在研究复变量的函数时特别有用。我们在第三、四章中讲过解析函数的微分与积分仍是解析函数，并且认为解析函数都定义在复平面 $\mathbb C$ 上的某个区域（甚至就是 $\mathbb C$ 本身）中，$\mathbb C$ 就是 $\mathbb R^2$。所以我们是在未扩充的 $\mathbb R^2$ 上讨论解析函数的。如果有某个函数 $f(z)$，令 $z=1/t$，而得到函数 $f(z)=f(1/t)=\varphi(t)$，它在 $0<|t|<1/R$ 中解析，如果 $\varphi(t)$ 一直到 $t=0$ 都是解析的，我们就说 $f(z)$ 在 $z=\infty$ 处也是解析的。这种观点可以说是在未扩充的 $\mathbb R^2$ 上把解析函数 $f(z)$ 解析延拓到 $\mathbb R^2\cup\{\infty\}$ 上。直到现在我们都是在未扩充的 $\mathbb R^n$ 讲微积分的。现在许多可采用另一种观点，即在扩充了的 $\mathbb R^2$，即 $\mathbb R^2\cup\{\infty\}$ 上补充这个函数，而把无穷远真正看作一个点。但这个观点与第一章绪论中讲的“实无穷”还有区别。如果看到，$\mathbb R^2\cup\{\infty\}$ 其实就是 $S^2$，而北极 $N\in S^2$ 是其上一个“普通的”点，这是任何人都能接受的。所以把 $\infty$ 看成扩充的 $\mathbb R^2$ 上一个“普通的”点也是十分自然的事。我们说这与“实无穷”还有区别是因为讨论“实无穷”时，我们甚至不需要找一个北极之类的对象来与之相应。但是这里究竟会遇到 $t\to0$ 时 $z\to\infty$，这个 $\infty$ 究竟不是通常的数。解决这个问题的办法是：不要每求一下 $z$ 就在整个 $S^2$ 上讨论函数而把 $S^2$ 分成几个部分：一部分不包括北极 $N$，例如就是 $\mathbb R^2$ 中的 $|z|<R$，在这一部分里讨论函数时，我们取 $z$ 为自变量，或者说以 $z$ 为“局部坐标”；另一部分则包括北极 $N$，也就是包括“扩充的”$\mathbb R^2$ 的“无穷远”点，例如 $|z|>R_1$（$R_1$ 与 $R$ 不一定相同而可取 $R_1<R$），这时我们就换一个新的局部坐标 $t$，而在 $|t|<1/R_1$（注意，包括 $t=0$，即 $z=\infty$）中讨论函数 $f(1/t)=\varphi(t)$。这两个部分可能相重，例如在 $R_1<|z|<R$ 中，我们则允许采用任一个局部坐标，而两种局部坐标 $t$ 与 $z$ 之间有一个“转换关系”（transition relation）$z=1/t$。这个十分自然的观点在现代数学中起了极大的作用，而拓扑的观点是其关键。下一章我们将要详细地展开它。至于在复变量函数的研究中，也就是在复分析中，在扩充的复平面上讨论解析函数也是极富成果的。这个扩充的复平面就称为黎曼球，这是黎曼对整个数学的伟大贡献之一。

可能使读者感到意外的是，$\mathbb R^2$ 不仅有这一种扩充的方法。为了阐释这一点，我们从一个表面上看来全无关系，甚至显得有些不自然的问题开始。不过我们声明一点，由于避免篇幅过长，以及避免还要更多的准备知识，下面关于射影平面的讨论将更多地依靠直观，而不甚严格。但是这里讲的结果都是正确的，都有严格的解释和证明。

考虑 $\mathbb R^3$ 中过原点之所有直线之集合 $X$，我们把每一条这样的直线看成一个点，并用下面的方法对这些点的集合赋以拓扑，并讨论这样得到的拓扑空间。设一条这样的直线有方向数 $(\alpha_1,\alpha_2,\alpha_3)$，这些 $\alpha_i$ 不能同时为0：$\sum_{i=1}^3\alpha_i^2\ne0$。于是，任取一个非零实数 $\lambda(\ne0)$，$(\lambda\alpha_1,\lambda\alpha_2,\lambda\alpha_3)$ 仍是同一条直线的方向数，所以我们不能说每一条这样的直线 $l$ 必与唯一一组 $(\alpha_1,\alpha_2,\alpha_3)$（其中 $\sum \alpha_i^2\ne0$）一一对应，而应该在这些三元组之间建立一个等价关系：
\[
 (\alpha_1,\alpha_2,\alpha_3)\sim(\beta_1,\beta_2,\beta_3)
 \iff
 (\alpha_1,\alpha_2,\alpha_3)=\lambda(\beta_1,\beta_2,\beta_3).
\]
这里 $\alpha=(\alpha_1,\alpha_2,\alpha_3)$，$\beta=(\beta_1,\beta_2,\beta_3)$ 都属于 $\mathbb R^3\setminus\{0\}$，这些 $\alpha,\beta$ 称为这条直线的齐次坐标（homogeneous coordinates）——其实并不是严格意义下的坐标，因为直线 $l$ 与方向数 $\alpha$ 并非一一对应，它只与 $\alpha$ 在关系 $\sim$ 下的等价类一一对应。这些等价类之集合我们记作 $\mathbb R^3/\sim$，称为 $\mathbb R^3$ 对于 $\sim$ 的商空间（quotient space）。商空间的概念是数学中的基本概念之一，其实我们过去已见过很多。例如周期运动，一个质点沿单位圆周作旋转，其所走的路程当然构成整个实数域 $\mathbb R$，但是每转一圈（正向或反向），这个质点都回到原来的位置，所以如果不问它已转了多少圈，而只问它现在在圆周上的位置，则运动路程 $s$ 和 $s+2k\pi,k\in\mathbb Z$ 其实代表同一位置——或称相位（phase）。所以我们用 $s\sim s+2k\pi$ 表示这个情况。这个质点走过的路程 $s\in\mathbb R$，但其相位之集合则为 $\mathbb R/\sim=\mathbb R/2\pi\mathbb Z$。相位空间是路程空间关于旋转整数周的商空间。现在回到原来的问题。在 $X$ 的每个元，即一条过 $O$ 的直线 $l$ 的等价类中选一个代表元。为此，在 $\mathbb R^3$ 中以 $O$ 为心作单位球面 $S^2:x_1^2+x_2^2+x_3^2=1$，于是每一个 $l$ 均与它交于两点，设其中一是 $x$（用交点在 $\mathbb R^3$ 中的直角坐标来表示），则另一交点是 $-x$，这两点位于一条直径的两端，称为对径点（antipodal points）。我们可以有两个说法，一是在这两个对径点中任取一点为其代表元，例如取位于下半球面 $S_-^2$ 上的一点代表 $l$；但这时位于赤道平面上的一对对径点，例如图6-3-2中的 $a$ 与 $a'$ 又如何？我们也可以规定取前面的一点，或者我们把 $S^2$ 当成地球表面，认为 $NA$ 为始初子午线，并由它开始计算经度，并且取经度在 $[0,\pi)$ 中的点为代表元，这样 $X$ 中的每一条直线均与 $S_-^2\cup\{\text{前半条赤道}\}$ 上的点一一对应，我们记 $S_-^2\cup\{\text{前半条赤道}\}$ 为 $\widetilde S_-^2$。则 $X$ 与 $\widetilde S_-^2$ 作为集合是一一对应的。这样的说法固然直观，但是太多人为而不自然的痕迹。尤其是，在应用它作具体计算时很不方便，因此我们又想到利用商空间的技巧，即认为 $x$ 与 $-x$ 互相等价：$x\sim(-x)$，而 $X$ 与 $S^2/\{x\sim(-x)\}$ 一一对应。对此有一个重要的说明：第一种说法的好处仍然在于其直观性，回到上面讲的转动的例子，起点位于 $S=0$ 与 $S=2\pi$，其实是一个点，就是说，如果把线段 $[0,2\pi]$ 的两端“粘合”起来，它就是圆周，$S=0$ 与 $S=2\pi$ 是等价的；具有相同相位，我们就把它们“粘合”在一起。把这个想法用于我们的例子：首先把 $S_+^2$ 去掉，只留下半球面 $S_-^2$，然后再把下半球面的边缘——赤道——上所有的对径点，例如图6-3-2中的 $a$ 与 $a'$，粘在一起，这样就可以得到 $\widetilde S_-^2$。而仍有 $X$ 与 $\widetilde S_-^2$ 作为集合一一对应。到这里读者会问，具体怎样粘法？举一个日常生活中的例子，如果你掏出了一张圆形的饺子皮，为了准备包馅，先把它拉成碗形，我们说这就是 $S_-^2$，可是最后怎样按上述要点捏成一个饺子呢？——要求饺子皮边上的每一对对径点要捏在一起。我们直觉的经验是：这是捏不成的。您可以做一个黎曼球那样的“饺子”，但做不成 $\widetilde S_-^2$ 饺子。为什么做不成？数学上可以证明，只有在4维空间才可以做成。这里我们回避了许多重要的数学问题：什么是商空间？它有什么性质？什么叫在4维空间才能做成？这些属于拓扑学的范围。但是为了看出它的困难所在，我们不妨再进一步，如果只是部分地把对径点捏到一起，例如已把 $S_-^2$ 边缘上用虚线画的两段粘在一起了：$a$ 点粘住了 $a'$，$b$ 点粘在 $b'$ 上，这就成为一个什么？把“饺子皮”$S_-^2$ 拉成一条长方形，把它的右端 $a'b'$ 转 $180^\circ$ 后粘到左边的 $ab$ 上就得到下面那个著名的曲面——莫比乌斯带（Möbius strip）（图6-3-3）。如果上面我们不把 $a'b'$ 转 $180^\circ$，而允许 $a'$ 粘到 $b$ 点，$b'$ 粘到 $a$ 点，就会得到一个柱面，它有上、下两个边，是两条曲线，可是莫比乌斯带的边只是一条曲线。请问，怎样把余下的部分 $ab'$ 与 $a'b$ 按我们的要求粘在一起呢？

\begin{figure}[htbp]
\centering
\begin{tikzpicture}[scale=.83, bella figure]
  % sphere with identified antipodal points
  \begin{scope}
    \draw[thick] (0,0) circle (2.05);
    \draw[thick] (-3.3,-1.2)--(3.3,-1.2)--(2.3,-.25)--(-4.0,-.25)--cycle;
    \draw[dashed] (-2.02,.05) arc[start angle=180,end angle=360,x radius=2.02,y radius=.55];
    \draw[thick] (2.02,.05) arc[start angle=0,end angle=180,x radius=2.02,y radius=.55];
    \draw[thick] (-1.3,.72)--(1.35,-.72);
    \fill (-1.05,.57) circle (1.5pt) node[left] {$a$};
    \fill (1.05,-.57) circle (1.5pt) node[right] {$a'$};
    \fill (-.85,-.50) circle (1.5pt) node[left] {$b$};
    \fill (.85,.50) circle (1.5pt) node[right] {$b'$};
    \node[above] at (0,2.05) {$N$};
  \end{scope}
  % rectangle
  \begin{scope}[xshift=7.1cm,yshift=.6cm]
    \draw[thick] (0,0) rectangle (4.3,1.0);
    \node[left] at (0,1) {$a$}; \node[left] at (0,0) {$b$};
    \node[right] at (4.3,1) {$b'$}; \node[right] at (4.3,0) {$a'$};
  \end{scope}
  % Mobius strip stylized
  \begin{scope}[xshift=7.7cm,yshift=-2.0cm]
    \draw[thick] plot[smooth cycle,tension=.7] coordinates {(-.3,0) (1.4,.9) (3.5,.2) (2.3,-1.0) (.6,-.6)};
    \draw[thick] plot[smooth cycle,tension=.7] coordinates {(.15,.08) (1.5,.55) (3.05,.12) (2.15,-.55) (.75,-.38)};
    \draw[thick] (1.25,.58)--(2.35,-.62);
  \end{scope}
\end{tikzpicture}
\caption{图6-3-2 与图6-3-3}
\end{figure}

回忆一下，我们在第二章中介绍了黎曼曲面，那也是一种无法在我们朝夕生活于其中的 $\mathbb R^3$ 中实现的。数学中出现了这样一些研究对象，并不只是由于想把微积分的基本概念弄得更明白。例如黎曼提出黎曼曲面是与他研究所谓阿贝尔积分密切相关的。更进一步，他终生思考的问题是，我们生活于其中的空间的物理本性究竟是什么？对于空间的研究与物理学的许多重大问题都相联系。拓扑学的创始者之一庞加莱研究拓扑学开始是由于天体力学中的所谓 $n$ 体问题，于是非常自然地出现了种种奇特的空间结构。正是为了研究这些空间结构，才进一步感到必须把一些最基本的概念弄清楚。这已经是20世纪初年的事了。

现在再回到上述关于 $X$ 与 $S^2$ 的研究。我们可以再提出一个与它们同胚的空间。在 $S^2$ 的南极 $S$\jiaokan{此处底本作“南极 $N$”，但按本节图6-3-1的记号及坐标 $x_3=-1$，南极应为 $S(0,0,-1)$，故改正。} 处作一个切平面 $x_3=-1$，在这个平面与 $S_-^2$ 间先建立一个作为集合的一一对应：若在 $S^2$ 上有一点 $P$，联结 $OP$ 交 $x_3=-1$ 于 $P_1$（见图6-3-4），反之，若在 $x_3=-1$ 上有一点 $P_1$，联结 $OP_1$ 交 $S_-^2$ 于 $P$。这样我们令 $P$ 与 $P_1$ 对应即得 $S_-^2$ 与 $x_3=-1$ 的对应。注意，这个对应还没有拓展到 $\widetilde S_-^2$ 上，因为联结 $O$ 与赤道上一点例如 $a$ 点，将得到一条与 $x_3=-1$ 平行的直线，它与 $x_3=-1$ 不会在任意有限远处相交。所以，若想把这个一一对应拓展到整个 $\widetilde S_-^2$ 上，就需要在此平面上添加一些理想元素，就如同把复平面变成扩充复平面或增加一个无穷远点 $z=\infty$，使它成为黎曼球一样。但是我们不能简单地在 $x_3=-1$ 上再添加一个点。因为在扩充复平面时，我们的目的是要把复变量的解析函数推广到 $z=\infty$ 附近，为此我们要利用 $z=1/t$ 以保持 $f(z)$ 的解析性质，而在交换 $z=1/t$ 之下，扩充的复平面上只许可有一个无穷远点 $z=\infty$ 与 $t=0$ 相应。现在我们拓展 $x_3=-1$ 的目的是什么呢？我们希望拓展后的平面尽可能多地保持欧氏平面的性质。在欧氏平面上有两个互相对偶的性质：(i)经过两个不同点可以作一条且仅有一条直线；(ii)两条直线，如果不平行，则必有一个且仅有一个有限远处的交点。请注意，第二个提法就是著名的平行线公设。这个公设在欧几里得《几何原本》中表述如下：“若一直线落在两直线上，所构成的同旁内角和小于两直角，那么把两直线无限延长，它们将在同旁内角和小于两直角的一侧相交。”直到1795年，普赖菲尔（John Playfair）才把这个公设改述为我们熟悉的形式：过直线外一点可以作一条，且仅可作一条与此直线平行的直线。所谓平行直线（《几何原本》卷I之定义23）是说，即“同一平面内的两直线，向两方无限延伸，而在两个方向上彼此不相交的直线。”可见，欧几里得是非常谨慎地避免说到“在无穷远处相交”一类的话，因为什么是“无穷远处”在数学上是没有定义的。现在我们扩充平面 $x_3=-1$ 想达到的目的就是把上述(i),(ii)两条加以推广，使(i)成为：(i')经过两个不同点（不论是有限远点还是无穷远点）都可以作一条且仅有一条直线；(ii)则推广成：(ii')两条不同直线必有一个且仅有一个交点（如果它们不平行，这个交点就是有限远点，如果平行，这个交点就是无穷远点）。这样一来，“平行线必在无穷远处相交”这个原本含义模糊的命题就变得完全确切清楚了。一个平面，这样添加上无穷远点以后，称为射影平面，记作 $\mathbb{RP}^2$；射影平面上的几何学称为射影几何学。我们立刻就会看到，射影平面与黎曼球面是截然不同的数学对象。由此可见，为了达到我们的目的，首先要说明什么是 $x_3=-1$ 上的直线。

如果在 $x_3=-1$ 上已经有了一条我们熟知的有限远处的直线（或线段）$l$（或 $P_1Q_1$），则 $OP_1,OQ_1$ 决定 $\mathbb R^3$ 中的一个平面，它与半球面 $S_-^2$ 沿一大圆弧 $aPQb$ 相交，大圆弧（劣弧）$PQ$ 对应于 $P_1Q_1$。在 $x_3=-1$ 上我们遇到了什么是无穷远点的困难，而是在球面 $S^2$ 以及 $S_-^2$ 上却没有这个困难。因此，我们把直观而朴素的直线“定义”推广为：$x_3=-1$ 上，大圆劣弧的像就是直线，大圆由通过球心的平面决定，大圆就是这个平面与球面之交。图6-3-4上我们画了一条通过有限远点 $P_1$ 的直线 $l$，它决定了相应大圆的平面，即 $OPQ$ 平面，它与赤道平面交于直径 $ab$。如上所说，$a\sim b$，而且我们以 $b$ 为代表元（$b\in\widetilde S_-^2$），$b$ 既然在大圆上，$b$ 之像就应该在 $l$ 上。在相应于 $l$ 的大圆弧 $(a,b)$（$a$ 不算弧上的点，但是 $b$ 算）上，只有这一点在赤道平面上，它的像不存在有限远点。因此，我们定义 $b$ 的像即 $l$ 上的无穷远点。这与我们的直观是很符合的，直线 $Ob$ 与平面 $x_3=-1$ 平行，所以它与此平面没有交点。但是从图6-3-4上看，由点 $P_1$ 沿 $l$ 向 $Q_1$ 移动时，球面上的相应点向 $b$ 移动，因此趋向无穷远。但由 $Q_1$ 向 $P_1$ 移动时，球面上相应点也趋向赤道平面，所以在 $l$ 上也趋向无穷远点，那么是否在 $l$ 上有两个无穷远点呢？没有，因为前一种情况下球面上的点趋向 $a$，而上面已经说过，$a$ 与 $b$ 是粘合在一起的；$a\sim b$，所以不论是向 $Q_1$ 移动还是向 $P_1$ 移动，都会趋向 $l$ 上的同一个无穷远点。总之我们看到，在 $x_3=-1$，即在 $\mathbb{RP}^2$ 上，每条通过至少一个有限远点 $P$ 的直线 $l$ 上均有恰好一个无穷远点，即大圆与赤道交点 $b$ 的像。上面我们加上了一点限制：“通过至少一个有限远点 $P$ 的直线 $l$”。

\begin{figure}[htbp]
\centering
\begin{tikzpicture}[scale=.92, bella figure]
  \draw[thick] (0,0) circle (2.0);
  \draw[dashed] (-2,0) arc[start angle=180,end angle=360,x radius=2,y radius=.55];
  \draw[thick] (2,0) arc[start angle=0,end angle=180,x radius=2,y radius=.55];
  \draw[thick] (-3,-2.0)--(3,-2.0)--(2.1,-1.15)--(-3.8,-1.15)--cycle;
  \coordinate (O) at (0,0); \coordinate (P) at (-.85,-1.25); \coordinate (Q) at (1.15,-1.0);
  \coordinate (P1) at (-1.45,-2.0); \coordinate (Q1) at (2.05,-2.0);
  \draw[thick] (O)--(P1); \draw[thick] (O)--(Q1);
  \draw[thick] (P1)--(Q1);
  \fill (O) circle (1.2pt) node[above] {$O$};
  \fill (P) circle (1.2pt) node[left] {$P$}; \fill (Q) circle (1.2pt) node[right] {$Q$};
  \fill (P1) circle (1.2pt) node[below left] {$P_1$}; \fill (Q1) circle (1.2pt) node[below right] {$Q_1$};
  \node[below] at (.3,-2.0) {$l$};
  \node[left] at (-2.0,0) {$a$}; \node[right] at (2.0,0) {$b$};
\end{tikzpicture}
\caption{图6-3-4}
\end{figure}

那么有没有不通过任何有限远点的直线呢？注意到，每一条直线都相当于一个大圆（其上的对径点视为一点），则还有一个大圆的像不在有限平面上，这就是赤道平面上的大圆（或如我们前面说的，只是这个大圆上经度在 $[0,\pi)$ 上的一段）。所以我们在平面上再引入一条理想的直线，称为无穷远直线。因为其上的点都是某一直线 $l$ 上的无穷远点，而没有有限远点。这样一来，通过两个不同无穷远点也有一条直线，即赤道平面上的大圆，这条直线就是唯一的不通过任何有限远点的直线。总之我们看到，为了保证 $\widetilde S_-^2$ 与平面 $x_3=-1$ 的点一一对应，就应在平面上增添一些“理想元素”，每条有限直线（即通过至少一个有限远点的直线）上加上一个（且仅加一个）无穷远点，以对应于赤道平面上之大圆上的点；整个平面应加上一条无穷远直线（它不通过任何有限远点），以对应于赤道平面上之大圆。添加了这些理想元素的平面称为射影平面 $\mathbb{RP}^2$。

我们可以多说几句。每条有限直线加上一个无穷远点，这在直观上不难理解，可是这些无穷远点的集合何以要称为一条直线，而不简单地就说是“无穷远点轨迹”之类？一方面因为这个轨迹是赤道平面上之大圆的像，而大圆之像总应该说是直线。还有另一个理由：我们一开始就讨论 $\mathbb R^3$ 中通过原点的直线，并且把每一条这样的直线均给以“齐次坐标”，即其方向数 $(a_1,a_2,a_3)$，这三个分量不许同时为0。其实在射影几何中，平面上每个点 $(x,y)$ 也都有一个齐次坐标：令 $x=x_1/x_3,y=x_2/x_3$，$(x_1,x_2,x_3)$ 就称为此点的齐次坐标，我们也规定这三个分量不得同时为0。我们知道，平面上的直线方程是
\[
 a_1x+a_2y+a_3=0,
\]
其中 $a_1,a_2$ 不许同时为0，但引入齐次坐标后，它成为
\begin{equation}
 \sum_{i=1}^3 a_i x_i=a_1x_1+a_2x_2+a_3x_3=0.
 \tag{4}
\end{equation}
因为 $(x_1,x_2,x_3)$ 地位平等——即地位对称——我们不应对哪一个有所偏爱，所以这时应加的限制就只能是 $a_1,a_2,a_3$ 不得同时为0（否则(4)式成为平凡的恒等式）。若 $a_1=a_2=0,a_3\ne0$，则(4)成为
\begin{equation}
 x_3=0,
 \tag{5}
\end{equation}
它仍是一条直线（一次方程），称为无穷远直线。同样，一个点 $(x,y)$ 即 $(x_1,x_2,x_3)$ 中，本应 $x_3\ne0$，但由于这三个分量现在是对称的，我们不应“歧视”$x_3$，不应该不允许它为0，而应规定 $x_1,x_2,x_3$ 不能同时为0（对方向数也有类似规定，若 $a_1=a_2=a_3=0$，就好比只有一个原点，作不出任何直线了）。我们规定，凡齐次坐标为 $(x_1,x_2,0)$ 之点称为无穷远点，当然它们的集合是无穷远直线(5)。从这一段讨论可以看到，无穷远点、无穷远直线与有限点、有限直线比较，没有任何本质上的特殊之处。这正是射影几何学的精义之一。

有了这些讨论，就可以看到欧氏平面上那两个对偶的性质现在成为

\noindent(i$'$) 过任意两个不同点必有唯一一直线。

当这两点都是有限远点（如图6-3-4上的 $P_1,Q_1$），这直线就是 $S_-^2$ 上的大圆弧 $PQ$；若一点是有限远点 $P$，一点是无穷远点 $b$，就是过 $P$ 与 $b$ 的大圆弧，它的像即过 $P_1$ 而与直径 $ab$ 平行的直线；若两点都是无穷远点，这直线就是无穷远直线。

\noindent(ii$'$) 任意两条不同直线必有唯一交点。

这是因为两个大圆必有交点，也是因为两个不同的方程(4)
\[
 a_i x_1+b_i x_2+c_i x_3=0,\qquad
 a_i^2+b_i^2+c_i^2\ne0,\qquad i=1,2,
\]
且 $(a_1,b_1,c_1)$ 与 $(a_2,b_2,c_2)$ 不等价，必有非0的公共解。如果这个交点不在赤道平面上而是有限远点，这就是初等几何中讲的两直线相交于某有限远点。如果这个交点是赤道平面上某点例如 $b$ 点，则从初等几何看来，这两条直线必平行于直径 $ab$，而有相同的无穷远点 $b$ 为其公共点。这一点读者很容易用立体几何的简单知识证明。

那么普赖菲尔的平行公设怎么会说呢？如果平行线理解为不相交的直线，那么这个公设错了，因为在我们的射影平面中没有两条不相交的直线，因此不存在所谓平行线。若我们把公设理解为过 $l$ 外一点必有一条直线与 $l$ 相交于其无穷远点 $b$，那么它就是(i$'$)。总之，在射影几何中没有给“平行性”留下特殊的地位，或者干脆就说平行性不是射影几何的概念。

倒是有一点应该强调，(i$'$)与(ii$'$)是对偶的：把“点”与“直线”互换一下，(i$'$)就变成(ii$'$)，(ii$'$)就变成(i$'$)。射影几何中许多极基本的定理都是对偶的，只要证明了其一，把其中的点与直线互换就得到另一个对偶的定理。初等几何中的许多定理都有这样的性质。而在现代数学中点 $x$ 与方向数所代表的向量 $\xi$ 的对偶关系，起了十分重要的作用。我们在第五章中还联系着量子力学讨论过这种对偶关系。这种对偶关系和射影几何的对偶关系本质上是相通的。

我们似乎离题太远，这是有原因的。射影几何理论源于透视法，在建筑学与绘画中都是基本的。在19世纪中，射影几何成了数学发展的主流的一个突出部分。当然，时过境迁，现在可以说是风光不再了。然而，它的许多基本概念都成了数学中不可少的财富。由于它与莫比乌斯带等等有密切关系，而莫比乌斯带在许多问题上又是典型，如它的单侧性、不可定向性（在下一章我们将看到这些概念在微积分中是多么重要），所以，对射影几何的一些基本概念还是知道一些为好。因为现在大学的数学课程一般都不再讲有关内容了，我们借此机会在这里多占用一些篇幅讲了以上的内容。下面还是回到拓扑学的问题。

我们上面实际上是讲了三个集合，即 $X$（$\mathbb R^3$ 中过原点之直线的集合）、$\widetilde S_-^2$ 以及 $\mathbb{RP}^2$ 之间有一一对应关系。但是重要的是如何在其上建立拓扑，从而使得这些集合之间的一一对应成为拓扑空间的同胚。这里最容易的从 $S^2$ 开始，$S^2$ 上如何定义开集似乎直观上没有什么困难。我们也可以一句以蔽之：让 $S^2$ 赋有 $\mathbb R^3$ 的子空间拓扑，即 $U\subset S^2$ 为一开集，当且仅当有一个 $\mathbb R^3$ 的开集 $V$，使 $U=V\cap S^2$。进一步我们来定义 $\widetilde S_-^2$ 上的拓扑，由于 $\widetilde S_-^2$ 是 $S^2$ 在一种等价关系 $\sim$ 下的商空间，我们说 $U\subset\widetilde S_-^2$ 是一个开集，意义如下：任找一个 $x\in U$，它在 $S^2$ 中所定义的等价类是 $S^2$ 的一个子集 $\{x,-x\}$，把这些子集（注意，不是看作一个等价类而是看成一个子集）并起来，得到一个 $V$（用一个很容易看懂的写法：$V=U\cup(-U)$），如果 $V$ 是 $S^2$ 的一个开集，我们就说 $U$ 是 $\widetilde S_-^2$ 的一个开集，这就是拓扑学中在商空间中定义商拓扑的方法。这件事看来简单，但有一些细微之处要注意，例如看图6-3-2，在 $A$ 点附近的 $U$，它在 $S^2$ 中的等价类之并分成两块，一块是 $S^2$ 上 $A$ 附近的 $U$，另一块是“下半球面”$S^2$ 上 $B$ 点附近的 $U$（图上的阴影区域）。如果 $A$ 点附近的 $U$ 是包括了赤道大圆弧的话，则 $B$ 附近的 $U$，也应包括赤道大圆弧。这两块并起来并非 $S^2$ 中的开集。但若在 $S^2$ 中取 $U\cup(-U)$，而且包括赤道大圆弧，即图6-3-2中分离的两块 $U$，它们一在 $A$ 点附近，一在 $B$ 点附近，则它在 $S^2$ 中的等价类之并合成两块，即图6-3-2上的两块连通的 $U\cup(-U)$，它们都跨过赤道大圆，而且都是 $S^2$ 上的开集。因此其并也是开集，所以 $\widetilde S_-^2$ 上分离的两块 $U$ 合在一起算是 $\widetilde S_-^2$ 的开集。有了 $\widetilde S_-^2$ 上的拓扑，即可定义 $\mathbb{RP}^2$ 与 $X$ 上的拓扑。我们来看前者，如果记图6-3-4\jiaokan{此处底本作“图6-3-3”，但由 $O,P,P_1$ 所示的投影关系实际见图6-3-4，故改正。}上由 $O$ 到 $P$ 到 $P_1$ 之映射为 $\varphi$，我们说 $U\subset\mathbb{RP}^2$ 为一开集，如果它是 $\widetilde S_-^2$ 上某开集 $V$ 的像的话：$U=\varphi(V)$。按前面的说法，这样在 $\mathbb{RP}^2$ 上定义拓扑仍是保证 $\varphi$ 为连续映射的最强的拓扑，但同时也是使 $\varphi^{-1}:\mathbb{RP}^2\to\widetilde S_-^2$ 为连续的最弱的拓扑。所以这是使 $\varphi$ 为同胚（即 $\varphi$ 与 $\varphi^{-1}$ 均为连续）的唯一的拓扑。对 $X$ 也可同样规定其拓扑。

总之，我们得到的不只是 $X,\widetilde S_-^2$ 与 $\mathbb{RP}^2$ 作为集合一一对应，而且是它们作为拓扑空间的同胚。这三个空间实际上是相同的，是射影平面的三个不同的表示方法。

把它与黎曼球比较就会感到，原来两个局部地看来完全一样的拓扑空间（二者局部地看来都是 $\mathbb R^2$ 的一部分），整体地看却完全不同。这样不由得提出一个问题：给出两个拓扑空间 $X$ 与 $Y$ 怎样看出它们是否同胚的呢？这是拓扑学的根本问题，也是无法完全解决的问题。

研究了关于拓扑空间及其映射的基本概念以后，读者可能会想到，是否应该继续讨论有关列极限等等问题了。很遗憾，不可以。通常凡讲到序列问题时，总是联系着度量空间的。关于序列的一些基本定理在一般拓扑空间中并不成立。当然，在微积分中有一些极限概念不能包括为序列极限的，例如在第四章中就讲到，黎曼积分的积分和的“极限”应该说就是上、下和的下、上确界更好。为了讨论这一类极限问题，以及把极限与收敛性问题放在拓扑空间来讨论，就需要引入“网”（net）的概念，这是拓扑学的另一重要概念，但在此割爱。

\noindent\textbf{3. 紧拓扑空间}\quad
现在回到拓扑空间的两个重要的特性：紧性与连通性。在这里我们将看到一般的拓扑空间与前面讲过特殊的拓扑空间——度量空间有多么大的差别。同时也把一些可以在上一节就开始讨论而没有展开的问题在这里继续下去。但在这之前，我们先介绍一个似乎不起眼而又十分重要的概念——可分离性。

\noindent\textbf{定义11}\quad
若 $(X,\tau)$ 为一拓扑空间，若 $X$ 中任意两个不同点 $x_1$ 与 $x_2$，均有其邻域 $U_i\ (i=1,2)$，使 $U_1\cap U_2=\varnothing$，这时 $X$ 称为豪斯多夫空间（Hausdorff space），或可分离空间。

这个概念似乎是完全平凡不足道的，它无非说 $X$ 中的两个不同点 $x_1$ 与 $x_2$ 总是可分离的。在度量空间中它必然成立（读者可以自己证明），但在一般的拓扑空间中它确实可能不成立。举一个平凡不足道的例子，若 $\tau$ 是平凡的拓扑，即只有 $\varnothing$ 与 $X$ 两个开集，这种不相交的邻域就找不到了。可是在数学中确实还有非平凡的拓扑空间也是非豪斯多夫的，这说明拓扑空间确实是一个极广泛的概念，我们有时不得不稍加限制使它更接近我们熟悉的度量空间。引入这个概念目的在此。附带提一句，引入 $(X,\tau)$ 的可分离性还是因为下一章将讨论的微分流形必须是可分离的。若不然，微积分中许多重要工具都会失效。

下面我们开始讨论紧性。对于一般的度量空间，§1中已经讲了的关于 $\mathbb R$ 的紧性的结果（§1，定理5）现在要作为定义来用。

\noindent\textbf{定义12}\quad
设 $(X,\tau)$ 为一拓扑空间，若 $X$ 之任一开覆盖 $\{U_\lambda\}$ 中必有一个有限子覆盖，则称 $(X,\tau)$ 为紧空间。若 $A\subset X$ 为一子集，且 $A$ 作为 $(X,\tau)$ 之子空间是紧的，则 $A$ 称为紧集。

\noindent\textbf{注}\quad
上面说 $A$ 作为子空间即指 $A$ 上赋有由 $X$ 诱导而得的子空间拓扑。所以当我们说 $\{U_\lambda\}$ 是 $A$ 的开覆盖时，$U_\lambda$ 就是子空间拓扑中的开集。亦即存在 $(X,\tau)$ 中的开集 $V_\lambda$，使 $U_\lambda=V_\lambda\cap A$。因此，在许多书中讲 $A$ 是紧集时说：“若有 $A$ 的开覆盖 $\{V_\lambda\}$，其中 $V_\lambda$ 是 $X$ 中之开集，则必有 $A$ 的有限子覆盖”，这与我们的讲法当然是一致的。

这个定义当然也适用于度量空间。

\noindent\textbf{定理6}\quad
若 $f:(X,\tau)\to(Y,\sigma)$ 为连续，$A\subset X$ 为紧集，则 $f(A)\subset Y$ 也是紧集。

\noindent\textbf{证}\quad
设 $\{V_\lambda\}$ 是 $f(A)$ 的开覆盖（这里的用法是指 $V_\lambda$ 是 $Y$ 中的开集），由于 $f$ 是连续的，故 $f^{-1}(V_\lambda)=U_\lambda$ 是 $(X,\tau)$ 的开集，而且 $\{U_\lambda\}$ 是 $A$ 的开覆盖，因为 $A$ 是紧集，故必有 $\{U_\lambda\}$ 的一个有限子集 $\{U_1,\ldots,U_N\}$ 使 $A\subset\bigcup_{k=1}^N U_k$，从而 $f(A)\subset\bigcup_{k=1}^N f(U_k)\subset\bigcup_{k=1}^N V_k$\jiaokan{底本此处写作 $\bigcup f(U_k)=\bigcup V_k$。由 $U_k=f^{-1}(V_k)$ 一般只能推出 $f(U_k)\subset V_k$；除非再有额外条件，等号并不成立。证明只需此包含关系，故正文改正。}，即是说 $\{V_1,\ldots,V_N\}$ 是 $f(A)$ 的有限子覆盖，从而 $f(A)$ 为紧。证毕。

\noindent\textbf{注}\quad
如果 $A=X$，这个定理仍成立。

紧集在连续映射下之像仍为紧，但 $(Y,\sigma)$ 中的紧集 $K$ 在 $(X,\tau)$ 中的原像则不一定为紧。如果对 $Y$ 之任一紧子集 $K$，$f^{-1}(K)\subset X$ 仍为紧，则 $f$ 称为适当的（proper）。

紧集在同胚映射下之像当然仍是紧的。这个简单的事实却告诉我们一个重要事实：若 $X=Y=\mathbb R$，$A=[a,b]\subset X$，$B=(a,b)\in Y$，则 $[a,b]$ 与 $(a,b)$ 决不可能同胚，否则 $(a,b)$ 将成为 $\mathbb R$ 的紧子集。但在 §1 中我们已经看到实数系的紧子集必为有界闭集，而 $(a,b)$ 不是有界闭集，所以它与 $[a,b]$ 不可能同胚。由此看来，开区间与闭区间的拓扑性质有很大的区别。到高维 $\mathbb R^n$ 时，情况也一样。所以在整个微积分理论中，我们时时都要注意是在开集还是在闭集上讨论问题，最简单如函数 $f(x)$ 的连续性，在闭区间上讨论时我们习惯把连续性了解为单侧连续性，只是最简单的例证。

紧性与闭性有密切的关系，我们有

\noindent\textbf{定理7}\quad
紧空间 $X$ 的闭子集 $A$ 必为紧集。

\noindent\textbf{证}\quad
设 $A$ 有一个开覆盖 $\{U_i\}$，因为 $A$ 为闭，所以 $A^c=X\setminus A$ 是 $X$ 的一个开子集，记作 $U^0$，于是 $\{U_i,U^0\}$ 构成 $X$ 的开覆盖，而存在一个有限子覆盖。不失一般性，我们可把 $U^0$ 也列入这个子覆盖中，而得有限子覆盖为 $\{U_1,\ldots,U_N,U^0\}$，所以
\[
 X=U^0\cup\left(\bigcup_{k=1}^N U_k\right).
\]
所以上式右方亦必覆盖 $A$。但 $U^0\cap A=\varnothing$，所以 $A\subset\bigcup_{k=1}^N U_k$，即是说从 $\{U_i\}$ 中可以找到 $A$ 的有限覆盖 $\{U_1,\ldots,U_N\}$，所以 $A$ 是紧集。

这个定理的逆则不太一样了：

\noindent\textbf{定理8}\quad
豪斯多夫空间的紧子集 $A$ 必为闭。

\noindent\textbf{证}\quad
为证 $A$ 是闭集，只需证明任一点 $x\in A^c=X\setminus A$ 必有一个与 $A$ 不相交的邻域 $V$ 即可。为此任取一点 $z\in A$，由于 $X$ 是豪斯多夫空间，故必有 $x$ 与 $z$ 的开邻域 $V_z$ 与 $U_z$ 互不相交：$V_z\cap U_z=\varnothing$，但 $\{U_z\}_{z\in A}$ 是 $A$ 的一个开覆盖，故必有一个有限子覆盖 $\{U_1,\ldots,U_N\}$ 使
\[
 A\subset\bigcup_{k=1}^N U_k=U.
\]
相应地则有 $x$ 的有限多个开邻域 $V_1,\ldots,V_N$ 使 $V_k$ 与 $U_k$ 不相交，记
\[
 V=\bigcap_{k=1}^N V_k,
\]
易见 $x\in V$，$A\subset U$ 而且 $V\cap U=\varnothing$。事实上，若有某点 $y\in V\cap U$ 则 $y$ 必属于某个 $U_j$，同样 $y\in V=\bigcap_{k=1}^N V_k\subset V_j$，于是 $V_j\cap U_j$ 中有至少一个元 $y$ 而非空集，这与 $V_j\cap U_j=\varnothing$ 相矛盾。

上面我们实际上证明了比定理所宣布的更多的结果：不但 $x$ 有一个邻域与 $A$ 不相交，而且还与 $A$ 的某个邻域 $U$ 不相交，而这又是下面结果的特例。

\noindent\textbf{定理9}\quad
设 $X$ 是豪斯多夫空间，$A_1,A_2$ 为 $X$ 之互不相交的紧子集，则必存在 $A_1,A_2$ 之邻域 $W_1$ 与 $W_2$，使 $W_1\cap W_2=\varnothing$。

\noindent\textbf{证}\quad
任取一点 $x\in A_1$，用定理8的证法，必可找到 $x$ 的邻域 $V_x$（即上面证明中的 $\bigcap_{k=1}^N V_k$）与 $A_2$ 的邻域 $U_x$（即上面的 $\bigcup_{k=1}^N U_k=U$），使得 $V_x\cap U_x=\varnothing$，而 $A_2\subset U_x$。但是 $\{V_x\}_{x\in A_1}$ 构成 $A_1$ 的开覆盖，所以可以由其中选出有限个 $V_x$ 来：$\{V_1,\ldots,V_M\}$，使得 $A_1\subset\bigcup_{k=1}^M V_k=W_1$；在 $\{U_x\}_{x\in A_1}$ 中取出相应的 $\{U_1,\ldots,U_M\}$，则 $A_2\subset\bigcap_{k=1}^M U_k=W_2$，$W_2$ 是 $A_2$ 的开邻域，而且仿照上面的证法，$W_1\cap W_2=\varnothing$。定理证毕。

这个定理告诉我们，一个空间若为豪斯多夫空间，则不但其不同点是互相分离的，而且其不相交的紧子集也是互相分离的。

至此，我们试图来刻画一个拓扑空间的全部紧子集。最简单的情况当然是 $\mathbb R^n$，我们可以证明，$\mathbb R^n$ 中的全部紧子集，即全部有界闭集。所以早期一点的微积分教材时常不讲紧性这个概念，而只讲其有界闭集如何如何。不言而喻，$\mathbb R^n$ 必是豪斯多夫空间。

\noindent\textbf{定理10}\quad
$\mathbb R^n$ 的紧子集，即 $\mathbb R^n$ 的有界闭集。

\noindent\textbf{证}\quad
我们只对 $n=1$ 来证明，先证 $\mathbb R^1$ 中的有界闭集 $A$ 为紧集，这在本质上就是 §1 定理6（海涅—博雷尔定理），不过那里讲的只是闭区间 $[a,b]$ 而不是一般的有界闭集 $A$，所以要稍作修改。因为 $A$ 是有界的，所以必有一个有界闭区间 $[m,M]$，使 $A\subset[m,M]$。令 $A^c=[m,M]\setminus A$，则 $A^c$ 是子空间拓扑中的开集。设有 $A$ 的开覆盖 $\{U_i\}$，则 $\{U_i,A^c\}$ 构成 $[m,M]$ 的开覆盖，而由 §1 定理6，必有一个有限子覆盖 $\{U_1,\ldots,U_N,A^c\}$，我们一定可以把 $A^c$ 取走，这不会有影响。于是 $\{U_1,\ldots,U_N\}$ 成为 $A$ 的覆盖，从而 $A$ 为紧。

再证另一方面，设 $A$ 为 $\mathbb R^1$ 的紧子集，今证 $A$ 是 $\mathbb R^1$ 的有界闭子集。因为 $\mathbb R^1$ 是豪斯多夫空间，由定理8知 $A$ 为闭集，余下的只要证明 $A$ 为有界的即可，作开区间 $I_n=(-n,n)$，则 $\{I_n\}$ 是 $A$ 的一个开覆盖（注意 $A$ 中的元都是数，而 $\pm\infty$ 不是数，所以 $A$ 中任一个 $x$ 必位于某个开区间 $I_{n_k}$ 中）。由 $A$ 之紧性，必有有限多个 $I_{n_1},\ldots,I_{n_N}$ 使 $A\subset\bigcup_{k=1}^N I_{n_k}$，令 $M=\max_{1\le k\le N}n_k<+\infty$，则 $A\subset[-M,M]$，从而 $A$ 有界。证毕。

至此，我们已成功地刻画了 $\mathbb R^n$ 中的全部紧子集。那么读者会问，能否类似地刻画度量空间中的全部紧子集，即全部有界闭集，至少是刻画巴拿赫空间或希尔伯特空间的全部紧子集即为有界闭集呢？不行。而且这是有深刻的理由的，所以我们暂时把它放在一边而来进一步讨论紧性。不过从现在起我们要特别注意一般的拓扑空间与度量空间的区别了。至少，所有度量空间都是豪斯多夫空间（请读者自己证明），而一般拓扑空间则不一定是。

先看定义在一紧集 $A$ 上的连续函数的性质，这个连续函数 $f(x),x\in A$，可以看成是由 $A\subset X$（$X$ 是一个拓扑空间）到 $\mathbb R$ 的连续映射。

\noindent\textbf{定理11}\quad
设 $(X,\tau)$ 是一个拓扑空间，$A\subset X$ 是 $X$ 的紧子集，$f(x):A\to\mathbb R$ 是 $A$ 上的连续函数，这时 $f(x)$ 必定是有界函数，而且可以达到其上、下确界。

\noindent\textbf{证}\quad
这个定理就是 §1 定理7的(i)，我们换一个更简单的证法，可以更清楚地看到它与度量性质没有关系，可以说它是一个纯粹紧性的证明。

由定理6，$f(A)$ 必为 $\mathbb R$ 之紧子集，由定理10又知 $f(A)$ 是有界闭集。既然 $f(A)$ 作为 $f$ 之值域是有界集，故 $f(x)$ 在 $A$ 上有上、下确界 $M=\sup f(x)$，$m=\inf f(x)$，又因 $f(A)$ 是闭集，其上、下确界必是 $f(A)$ 中之点，因此存在 $x_1,x_2\in A$ 使 $f(x_1)=M,f(x_2)=m$。定理证毕。

读者会问，§1 定理7之(ii)即关于一致连续性又如何？一致连续性是一个度量性质，因此(ii)在一般拓扑空间上的推广要利用一种介乎度量结构与拓扑结构中的概念，我们就不再讨论了。

我们在微积分教材中都会看到一个基本定理，即任一无穷集合至少有一个极限点仍在此空间中。若一拓扑空间有此性质，就说它具有魏尔斯特拉斯—波尔查诺性质。要注意，由定义6，极限点之定义只与邻域有关。它是一个一般的拓扑性质，而与序列概念无关，$x$ 是极限点即指 $x$ 之任一邻域中有集合 $A$ 中与 $x$ 不同的点，这与存在一个收敛于 $x$ 的子序列的概念是两回事，这一点必须十分注意。

\noindent\textbf{定理12}\quad
紧空间 $X$ 必具有魏尔斯特拉斯—波尔查诺性质。

\noindent\textbf{证}\quad
设 $S\subset X$ 没有极限点，今证 $S$ 必为有限集。

任取一点 $x\in X$，必可找到一个开邻域 $U_x$，使得
\begin{equation}
 U_x\cap S=
 \begin{cases}
 \varnothing, & \text{若 }x\notin S,\\
 \{x\}, & \text{若 }x\in S,
 \end{cases}
 \tag{6}
\end{equation}
否则的话，$x$ 将是 $S$ 的极限点。$\{U_x\}_{x\in X}$ 是 $X$ 的一个开覆盖，因为 $X$ 是紧的，故从 $\{U_x\}$ 中可以取出有限多个 $\{U_1,U_2,\ldots,U_N\}$ 来覆盖 $S$，但由(6)式，$U_k$ 中或者没有 $S$ 之元，或者只有一个元。所以 $S$ 必为有限集。

读者会问，这个定理之逆是否成立？如果成立，则我们可以就用魏尔斯特拉斯—波尔查诺性质作为紧性的定义，可以明确指出，这是不可能的。由它只能得到较弱的结论：从 $X$ 之任一可数开覆盖中可以找到有限子覆盖。因此，紧性的定义只能是定义12而不能用“无穷子集必有极限点”来代替，更不能用“任意序列均有收敛子序列”来代替。

下面我们转到度量空间中的紧性。为此，我们首先要给出一个技术性的结果。

\noindent\textbf{定理13（勒贝格引理）}\quad
设 $X$ 为一紧度量空间，$\{U_i\}$ 为其一个开覆盖，则存在一个常数 $\delta>0$（称为 $\{U_i\}$ 的勒贝格数），使 $X$ 之任一直径小于 $\delta$ 的集合必含于 $\{U_i\}$ 之某一元内。

\noindent\textbf{证}\quad
用 $\rho$ 表示 $X$ 中的度量，$B_r(P)$ 表示以 $P$ 为心、$r$ 为半径的球。一个集合 $U$ 之直径即
\[
 \sup_{x,y\in U}\rho(x,y).
\]
我们用反证法证明它。设定理结论不真，则必存在 $X$ 之子集序列 $A_n$，使所有 $A_n$ 均不含于 $\{U_i\}$ 之任一元 $U_i$ 内，但 $\operatorname{diam}A_n\to0$。在 $A_n$ 中选一点 $x_n$，于是得一无穷子集 $\{x_n\}$，由魏尔斯特拉斯—波尔查诺性质，这个子集有极限点 $P$，因为 $\{U_i\}$ 是 $X$ 的一个开覆盖，故必有一个 $U\in\{U_i\}$ 使 $P\in U$。今取 $\varepsilon>0$，使 $B_\varepsilon(P)\subset U$，再取 $N$ 足够大使得(i)$\operatorname{diam}A_N<\varepsilon/2$；(ii)$x_N\in B_{\varepsilon/2}(P)$，于是 $\rho(x_N,P)<\varepsilon/2$，而对 $A_N$ 之任何一点 $x$，$\rho(x,P)\le\rho(x,x_N)+\rho(x_N,P)<\varepsilon$。这表明 $A_N\subset U\in\{U_i\}$，这与 $A_N$ 不含于 $\{U_i\}$ 之任一元之内相矛盾。

\noindent\textbf{注}\quad
以上说 $\{x_n\}$ 是一无穷子集需要解释，因为一个集合中每个元只能计一次（而序列则不然，若说 $\{x_n\}$ 是一序列，即定义在自然数集上的函数，则同一个 $x_n$ 可以出现多次，例如 §2 讲到度量空间完备化时讲到的常驻序列），而这里的 $x_n$ 取自 $A_n$ 却是可以重复的。若 $x_{n_0}$ 重复了无穷多次，我们即可用 $x_{n_0}$ 代替极限点，也就取 $x_{n_0}$ 为 $x_N$；若 $x_n$ 只重复有限次，则只计一次还会留下无穷多个不同的 $x_n$，而上面的讨论可以顺利进行。

现在我们就来讨论度量空间的紧性，注意，它与一般拓扑空间是有区别的。

首先来看一致连续性，设 $X,Y$ 是两个度量空间，其度量分别是 $\rho$ 与 $\sigma$，我们说映射 $f:X\to Y$ 是一致连续的（或者准确一点说是关于 $(\rho,\sigma)$ 一致连续的）是指对任意 $\varepsilon>0$，必可找到常数 $\delta(\varepsilon)>0$，使当 $\rho(x_1,x_2)<\delta$ 时，$\sigma(f(x_1),f(x_2))<\varepsilon$，$\delta(\varepsilon)$ 与 $X$ 中取 $x_1,x_2$ 之取法无关。

\noindent\textbf{定理14}\quad
若 $X,Y$ 为度量空间，其度量为 $\rho$ 和 $\sigma$，若 $X$ 为紧，则一切连续映射 $f:X\to Y$ 必在 $X$ 上一致连续。

\noindent\textbf{证}\quad
因为 $f$ 已设为连续，故对任给的 $\varepsilon>0$ 以及 $x\in X$ 必可找到 $\delta(\varepsilon,x)>0$，使当 $\tilde x\in B_\delta(x)$ 时，$\sigma(f(\tilde x),f(x))<\varepsilon/2$。球 $\{B_\delta(x)\}$ 构成 $X$ 的开覆盖，因此有一个有限子覆盖 $\{B_{\delta_k}(x_k)\}$，$k=1,2,\ldots,N$。相应于这个覆盖，必有其勒贝格数 $\delta(\varepsilon)>0$，它与 $x_k$ 无关。今设有任意两点 $x'$ 与 $x''$ 适合 $\rho(x',x'')<\delta$，这时 $x'$ 与 $x''$ 必同属于同一个 $B_{\delta_k}(x_k)$，因此
\[
 \sigma(f(x'),f(x''))\le \sigma(f(x'),f(x_k))+\sigma(f(x_k),f(x''))
 <\frac{\varepsilon}{2}+\frac{\varepsilon}{2}=\varepsilon.
\]
定理证毕。

在许多微积分教本中曾以任一无穷序列 $\{x_n\}$ 均有收敛子序列作为紧性的定义，这样做明显是有问题的，因为前面我们一再提到，把有关极限的问题归结为序列问题一般说来适用于度量空间，但是有一个重要的例外，广义函数理论中的那些函数空间一般不是度量空间，但是在第四、五两章中讲到广义函数空间的拓扑性质时，总是使用序列。这是需要专门研究的，本书中我们不能去做这种研究，但是应该指出这样做是必需的，而且是可能的。在紧性概念问题上我们已经有了三个概念：以覆盖为基础的紧性定义；以魏尔斯特拉斯—波尔查诺性质为基础的紧性定义；现在又有了以子序列收敛为基础的紧性定义，更确切一点，有

\noindent\textbf{定义13}\quad
设 $X$ 为一度量空间，若 $X$ 之任一无穷序列 $\{x_n\}$ 都有收敛的子序列，即有子序列 $\{x_{n_k}\}$ 收敛于 $X$ 中某个元，则称 $X$ 为列紧空间（sequentially compact space）。

我们有以下的基本定理。

\noindent\textbf{定理15}\quad
若 $X$ 为一度量空间，则以下三个命题等价：

\noindent(i) $X$ 为紧；

\noindent(ii) $X$ 具有魏尔斯特拉斯—波尔查诺性质；

\noindent(iii) $X$ 为列紧。

\noindent\textbf{证}\quad
(i)$\Rightarrow$(ii)：即定理12应用于度量空间。

(ii)$\Rightarrow$(iii)：设 $\{x_n\}\subset X$ 为一无穷序列。定理13的注中我们强调了，一个序列 $\{x_n\}$ 与子集 $\{x_n\}$ 之区别在于前者允许有相等的 $x_n$，后者则不允许。今设 $\{x_n\}$ 中有无穷多个彼此相等的元，即有 $x_{n_1}=x_{n_2}=\cdots=x_0$，于是子序列 $\{x_{n_k}\}$ 是一个常驻序列，它自然收敛于 $x_0\in X$。若 $\{x_n\}$ 中没有无穷多个相等的元，则可以找到一个子集仍记为 $\{x_n\}$（即将相等的元只留一个以后删去其多余者），它是一个无穷子集。由魏尔斯特拉斯—波尔查诺性质，它必有一个极限点 $x_0\in X$，但 $x_0$ 之任一邻域中均可找到 $\{x_n\}$ 之至少一个与 $x_0$ 不同的（这是由极限点之定义6得知的）点。今任取这样的一个点 $x_{n_1}$，并令 $\rho(x_{n_1},x_0)=\varepsilon_1>0$，于是在 $B_{\varepsilon_1/2}(x_0)$ 中又可找到一点 $x_{n_2}$，使 $\rho(x_{n_2},x_0)=\varepsilon_2>0$，且 $\varepsilon_2<\varepsilon_1/2$，仿此可得一个子序列 $x_{n_k}$，使 $\rho(x_{n_k},x_0)=\varepsilon_k<2^{-(k-1)}\varepsilon_1$，而且 $\varepsilon_k>0$，很明显 $\lim_{k\to\infty}x_{n_k}=x_0$，而(iii)得证。

(iii)$\Rightarrow$(i)：设 $X$ 为列紧而 $\{U_i\}$ 是 $X$ 的一个开覆盖。由勒贝格引理（这个结果是对紧空间证明的，但是读者容易看到，定理13之证明对列紧空间也是有效的），这个开覆盖应有勒贝格数 $\varepsilon$。今取 $0<\delta<\varepsilon$ 并证明必有有限多个以 $\delta$ 为半径的球 $B_1,\ldots,B_N$，使 $\bigcup_{k=1}^N B_k\supset X$，这一点可用反证法证明。若它不对，则以某个 $x_1$ 为心，$\delta$ 为半径的球 $B_1=B_\delta(x_1)$ 不能覆盖 $X$，故有 $x_2\in X\setminus B_\delta(x_1)$，再作 $B_2=B_\delta(x_2)$，则 $B_1\cup B_2$ 仍不能覆盖 $X$，故必有 $x_3\in X\setminus(B_1\cup B_2)$，而 $\rho(x_1,x_3)\ge\delta,\rho(x_2,x_3)\ge\delta$，再作 $B_3=B_\delta(x_3)$，并找到 $x_4$。仿此，可以找到一个无穷序列 $\{x_n\}$，使 $\rho(x_n,x_m)\ge\delta,n\ne m$。这样一个无穷序列 $\{x_n\}$ 因为其中任何两个元距离不小于 $\delta>0$，而它不可能有收敛子序列，这与 $X$ 之列紧性矛盾。定理证毕。

于是我们看到在度量空间中这三个有关紧性的概念是等价的，而在非度量空间则不一定。这就是在一般拓扑空间中三个概念不能混淆的原因。

下面我们进而讨论一个在许多数学分支中都一再出现的有关紧性的问题，即给出一个空间或其中一个子集，弄清楚在什么条件下它们是紧的。因为我们现在主要关心的已不再是有关紧性的一般概念，我们将把注意力放在最常见的空间上，首先是赋范线性空间（包括内积空间）。这些空间都是有线性结构的，因此也就有维数，而我们很快就发现，维数在这里起了决定作用。

§1 中我们已看到，$\mathbb R^1$ 是局部紧的，即任一点都有一个邻域，其闭包为紧。例如 $\mathbb R^1$ 上一点 $x_0$ 的邻域就可以取为含此点的区间，其闭包就是闭区间 $[a,b]$，而 $x_0$ 为其内点，它当然是紧集——即作为 $\mathbb R^1$ 之子空间，在子空间拓扑下是紧空间。$\mathbb R^1$ 是线性空间，直观地看就是一条平直的直线，其各点的邻域本质上是一样的，或者用数学语言表述，其上的平移变换是一对一的仿射（但不是线性，因为在平移下原点不是不变的）同构与拓扑同胚。所以要证明 $x_0$ 的邻域闭包为紧只需证明 $x=0$ 的邻域闭包 $[-a,a],a>0$ 为紧。再以0为中心作一个相似变换，它也是线性同构又加上同胚，而且把 $[-a,a]$ 变成1维单位球体 $[-1,1]$，所以证明 $\mathbb R^1$ 为局部紧，只需证明其以0为心的单位球体为紧即可，对于 $\mathbb R^n$ 也是一样。

$\mathbb R^n$ 为局部紧但本身并不紧，这是非常重要而又自然的。这都是 $\mathbb R^1$ 这个空间的“均匀性”造成的。要证明 $\mathbb R^1$ 非紧我们建议用以下证法：以整数 $k$ 为心，$2/3$ 为半径作开区间
\[
 I_k=\left(k-\frac23,k+\frac23\right),\qquad k=0,\pm1,\pm2,\ldots
\]
很明显 $\{I_k\},k=0,\pm1,\pm2,\ldots$ 是 $\mathbb R^1$ 的开覆盖。我们绝不可能从其中找到有限的子覆盖，因为要覆盖 $\mathbb R^1$ 就得覆盖所有的整点 $x=k,k=0,\pm1,\pm2,\ldots$，而每个整点 $x=k$ 恰好只落在一个 $I_k$ 中。所以想要覆盖这些整点以至覆盖整个 $\mathbb R^1$，这些 $I_k$ 一个也不能少。这正是 $\mathbb R^1$ 的“均匀性”的表现。

同理 $\mathbb R^n$ 也都是局部紧而本身非紧。但转到无穷维情况就完全不同了，我们下面只是并不严格地举例以说明，严格证明还要多费一些口舌。设 $H$ 是一个实无穷维希尔伯特空间，于是有一个 o.n. 系 $\{e_1,e_2,\ldots,e_n,\ldots\}$（我们不问它是否可分，因为不去讨论其完备性），它们都位于单位球面 $\|x\|=1$ 上。单位球面当然是一个度量空间，因此我们想利用定理15，看一下 $\|x\|=1$ 是否列紧。它当然不是，因为序列 $\{e_n\}$ 中两个不相同的元之距离平方为
\[
 \|e_m-e_n\|^2=\|e_m\|^2-2(e_m,e_n)+\|e_n\|^2=2,
\]
而不可能有收敛子序列。如果要严格地讲，我们不妨证明地给出

\noindent\textbf{定理（F. 里斯）}\quad
赋范线性空间 $B$ 为局部紧的充分必要条件是 $B$ 为有限维空间。

由此可见，在无穷维空间中一个集合为紧，除了有界条件以外，一定要加上其它条件。循着这个思路可以得到许多在各数学分支中十分有用的结果。下面我们来介绍其中可能是最有名的一个：阿斯科里—阿尔泽拉定理（以下称为 A--A 定理），阿斯科里（G. Ascoli）和阿尔泽拉（C. Arzelà）都是意大利数学家，他们二人都发现了这个定理。这个定理讲的是有界闭区间 $I=[a,b]$ 上的连续函数空间 $C(I)$ 中的集合为紧的充分必要条件，因为 $C(I)$ 是一个度量空间，所以其中的紧与列紧是一致的。在讨论这个问题前，先作两点说明。

所谓 $M$ 列紧，即 $M$ 中任一序列 $\{\varphi_n(x)\}$ 必有 $C(I)$ 意义下的收敛子序列 $\{\varphi_{n_k}(x)\}$：$\lim_{k\to\infty}\varphi_{n_k}(x)=\varphi_0(x)$。$C(I)$ 意义下的收敛即一致收敛。因此，$\{\varphi_{n_k}(x)\}$ 之极限 $\varphi_0(x)$ 仍在 $C(I)$ 中，即仍为 $I$ 上的连续函数，但不能确定是否在 $M$ 中，而只能确定它在 $\overline M$ 中，所以实际上我们只能证明 $\overline M$ 为列紧。不过由定理8，豪斯多夫空间——$C(I)$ 也是——中紧子集必为闭，所以应该证明的只是 $\overline M$ 为列紧。若一个集合之闭包为紧，则称此集合为预紧（precompact）。所以我们应该证明的是 $M$ 为预紧。

A--A 定理中要用到等度连续（equicontinuous）概念：若 $M$ 为 $C(I)$ 中之无穷子集（有限子集之等度连续性概念没有意义）为等度连续的，是指对任一 $\varepsilon>0$，必有 $\delta(\varepsilon)>0$ 存在，使当 $|x_1-x_2|<\delta(\varepsilon)$ 时，不论 $x_1,x_2$ 在 $I$ 中何处，也不论 $\varphi(x)$ 是 $M$ 之哪一个元，都有 $|\varphi(x_1)-\varphi(x_2)|<\varepsilon$。这就是说 $\delta(\varepsilon)$ 不仅不依赖于 $x$（这是一致连续性的条件），而且不依赖于 $\varphi(x)$ 之选择。于是我们给出 A--A 定理如下：

\noindent\textbf{定理16（A--A 定理）}\quad
$C(I)$ 之子集 $M$ 为预紧的充分必要条件是 $M$ 中的函数有公共的上界且为等度连续。

\noindent\textbf{证}\quad
\textbf{必要性}\quad 设 $M$ 为预紧。我们先用一个概念即 $\varepsilon$ 网来表述紧性。若度量空间 $X$ 为紧，以任一点 $P$ 为心、$\varepsilon$ 为半径作一球 $B_\varepsilon(P)$，则 $\bigcup_{P\in X}B_\varepsilon(P)$ 覆盖 $X$，因此可以找到有限个这样的球 $B_\varepsilon(P_k),k=1,2,\ldots,N$ 使 $\bigcup_{k=1}^N B_\varepsilon(P_k)=X$。$\{P_1,\ldots,P_N\}$ 就称为一个 $\varepsilon$ 网，所以若一度量空间为紧，则对任意 $\varepsilon>0$ 均有有限 $\varepsilon$ 网存在（其逆亦真，这一点下面要用到）。把它应用于 $M$，并用它来证明 $M$ 中之函数均为等度连续的。我们构造 $M$ 的 $\varepsilon/3$ 网，$\{\varphi_1(x),\ldots,\varphi_N(x)\}\subset C(I)$，每个 $\varphi_k(x)$ 均为 $I$ 上之连续函数，故为一致连续。因此对 $\varepsilon/3$，可以找到 $\delta_k(\varepsilon)$，使当 $|x_1-x_2|<\delta_k(\varepsilon)$ 时
\[
 |\varphi_k(x_1)-\varphi_k(x_2)|<\frac{\varepsilon}{3}.
\]
在有限多个 $\delta_k(\varepsilon)$ 中取 $\delta(\varepsilon)=\min\{\delta_1(\varepsilon),\ldots,\delta_N(\varepsilon)\}$，显然 $\delta(\varepsilon)>0$，于是对任一 $\varphi(x)\in M$，可以找到 $\varepsilon/3$ 网中的 $\varphi_k(x)$，而得到，当 $|x_1-x_2|<\delta(\varepsilon)$ 时
\[
 |\varphi(x_1)-\varphi(x_2)|
 \le |\varphi(x_1)-\varphi_k(x_1)|+|\varphi_k(x_1)-\varphi_k(x_2)|+|\varphi_k(x_2)-\varphi(x_2)|
 <\varepsilon.
\]
因此 $M$ 是等度连续的。

$M$ 中之元有公共的界易证，因为对上述 $\varepsilon/3$ 网，每个 $|\varphi_k(x)|$ 在 $I$ 上均有上界 $C_k$，令 $C=\max\{C_1,\ldots,C_N\}$，则 $C<+\infty$，而 $|\varphi_k(x)|\le C$。于是对 $M$ 中任一 $\varphi(x)$ 仿上法有
\[
 |\varphi(x)|\le |\varphi_k(x)|+|\varphi(x)-\varphi_k(x)|\le C+1,
\]
可见 $C+1$ 即 $M$ 中所有函数绝对值的公共上界。

\textbf{充分性}\quad 充分性的证明基本思想是用有限维空间去逼近 $\overline M$，应该强调指出，这是在涉及紧性问题时的一个基本方法。

对于一个连续函数 $\varphi(x)\in\overline M$，将 $I$ 等分为
\[
 a=a_0<a_1<\cdots<a_n=b,\qquad a_k=a_0+kh,\quad h=\frac{b-a}{n}.
\]
于是得到 $\varphi(x)$ 曲线上的 $n+1$ 个点 $(a_i,\varphi(a_i))$，以它们为顶点作折线，记其方程为 $y=\varphi^{(n)}(x)$。（请注意，这里的 $\varphi^{(n)}(x)$ 是指具有 $n+1$ 个顶点的折线的方程右方，而不是 $\varphi(x)$ 的 $n$ 阶导数。）从图形上很容易看到，在子区间 $[a_i,a_{i+1}]$ 中若 $\varphi(a_i)\le\varphi(a_{i+1})$ 必有
\[
 \varphi(a_i)\le\varphi^{(n)}(x)\le\varphi(a_{i+1}),
\]
而当 $\varphi(a_i)\ge\varphi(a_{i+1})$ 时则有
\[
 \varphi(a_{i+1})\le\varphi^{(n)}(x)\le\varphi(a_i).
\]
总之，当 $x\in[a_i,a_{i+1}]$ 时
\begin{align}
 \varphi(x)-\varphi(a_{i+1})&\le\varphi(x)-\varphi^{(n)}(x)\le\varphi(x)-\varphi(a_i), \tag{7$_1$}\\
 \varphi(x)-\varphi(a_i)&\le\varphi(x)-\varphi^{(n)}(x)\le\varphi(x)-\varphi(a_{i+1}). \tag{7$_2$}
\end{align}
但是 $\varphi(x)$ 是一致连续的，上式中 $|x-a_i|,|x-a_{i+1}|\le1/n$，故当 $n$ 充分大时，(7$_1$),(7$_2$)左右两方之项绝对值均小于 $\varepsilon$，因此对 $x\in I$ 有
\begin{equation}
 \sup_I|\varphi(x)-\varphi^{(n)}(x)|<\varepsilon.
 \tag{8}
\end{equation}
而得到 $\overline M$ 的一个由折线组成的 $\varepsilon$ 网，记为 $N_1$。

这些折线还是有界的，因为对每条折线 $\varphi^{(n)}(x)$ 而言
\[
 |\varphi^{(n)}(x)|\le |\varphi(x)-\varphi^{(n)}(x)|+|\varphi(x)|\le\varepsilon+C\le C+1.
\]
这里 $\varphi(x)$ 是 $\overline M$ 中适合(8)式的曲线，而 $C$ 是 $M$ 中元素绝对值的公共上界。

当然，这个 $\varepsilon$ 网不是有限的，但每一条折线均由固定的分点 $x=a_i$ 处折线的纵坐标 $b_i$ 决定，所以每一条折线对应于 $\mathbb R^{n+1}$ 中的一点 $(b_0,b_1,\ldots,b_n)$。因为每个 $|b_i|$ 都有界，故对应于上面得到的 $\varepsilon$ 网的 $(b_0,b_1,\ldots,b_n)$ 之集合包含在 $\mathbb R^{n+1}$ 的一个有界闭盒中；这个闭盒是紧集，因而该系数集合有一个有限 $\varepsilon$ 网，记为 $N_2$\jiaokan{底本此处由“每个 $|b_i|$ 都有界”直接断言这些系数点本身构成 $\mathbb R^{n+1}$ 中的紧集。有限维中有界并不自动推出闭，故不能据此断言该集合本身紧；但它包含在一个有界闭盒中，而该闭盒紧，从而仍可得到所需的有限 $\varepsilon$ 网。}。这个新的有限 $\varepsilon$ 网 $N_2$ 之每一点对应于一条折线 $\psi^{(n)}(x)$，取 $N_2$ 中的 $\psi^{(n)}(x)$，即有
\[
 |\varphi(x)-\psi^{(n)}(x)|\le |\varphi(x)-\varphi^{(n)}(x)|+|\varphi^{(n)}(x)-\psi^{(n)}(x)|<2\varepsilon.
\]
所以 $N_2$ 是 $\overline M$ 的 $2\varepsilon$ 网，因为 $2\varepsilon$ 是任意数，我们不妨仍记为 $\varepsilon$ 网 $N_2$。

这个 $N_2$ 中之函数绝对值仍有公共上界。

在必要性部分中我们已证明了，度量空间 $X$ 中任一列紧集均有有限 $\varepsilon$ 网。现在要证明其逆：若这个度量空间 $X$ 还是完备的，而闭子集 $\overline M$ 又对任意 $\varepsilon$ 有有限的 $\varepsilon$ 网，则此闭子集 $\overline M$ 必为列紧的。这两个结论合起来是豪斯多夫证明的一个定理。我们要特别注意其充分性部分中假设了 $X$ 是完备空间。

现在证明其充分性部分，取一串 $\varepsilon_n\downarrow0$，再取 $\overline M$ 中一个序列
\[
 S_0=\{\varphi_1,\ldots,\varphi_N,\ldots\}.
\]
先作 $\overline M$ 之 $\varepsilon_1$ 网 $\{\psi_1^{(1)},\ldots,\psi_{N_1}^{(1)}\}$，以它们为心以 $\varepsilon_1$ 为半径作球 $B_{\varepsilon_1}(\psi_k^{(1)})$，必可覆盖 $\overline M$。因此至少有一个这样的球包含了 $\{\varphi_1,\ldots,\varphi_N,\ldots\}$ 之无穷多个元，记为
\[
 S_1=\{\varphi_1^{(1)},\ldots,\varphi_N^{(1)},\ldots\}.
\]
下面只看序列 $S_1$，再作 $\overline M$ 的有限 $\varepsilon_2$ 网以及相应的以 $\varepsilon_2$ 为半径的球，于是又可找到 $S_1$ 的一个子序列全含于此球中：
\[
 S_2=\{\varphi_1^{(2)},\ldots,\varphi_N^{(2)},\ldots\}.
\]
仿此可以作出 $S_2$ 之子序列 $S_3$，$S_3$ 的子序列 $S_4,\ldots$，而每一个子序列都含于半径为 $\varepsilon_k\ (\downarrow0)$ 的球中。取这些子序列的对角线序列
\[
 \{\varphi_1^{(1)},\ldots,\varphi_2^{(2)},\ldots,\varphi_N^{(N)},\ldots\},
\]
请读者自己证明它是一个柯西序列，而因 $X$ 为完备的，所以它有一个极限 $\varphi\in X$。又因 $\overline M$ 是一闭集，所以 $\varphi\in\overline M$。这样得知，在 $\overline M$ 中任一序列 $S_0$ 都有一个收敛子序列收敛于 $\overline M$ 之某个元 $\varphi$，这样 $\overline M$ 是列紧的也就是紧的。定理证毕。

与数学中许多重大的定理一样，A--A 定理在 $L^p(I)$ 中也有由里斯证明了的相应的结果：设 $M\subset L^p(I)$，使得

\noindent(i) 存在一个常数 $K$，使得一切 $\varphi(x)\in M$ 均有
\[
 \int_I |\varphi(x)|^p\,dx\le K^p;
\]

\noindent(ii) 对任意 $\varepsilon>0$ 均有一个 $\delta(\varepsilon)>0$ 存在，使得一切 $\varphi(x)\in M$，只要 $|h|<\delta$ 就有
\begin{equation}
 \int_I |\varphi(x+h)-\varphi(x)|^p\,dx\le\varepsilon^p,
 \tag{9}
\end{equation}
这时，$\overline M$ 必为 $L^p(I)$ 中的紧集。

很容易看出，(9)就是 $L^p$ 意义下的等度连续性。它在上面的讨论中起了关键作用，因此我们自然要问，它在什么条件下可以成立？我们仍以 $C(I)$ 为例说明它。若 $\varphi(x)$ 不仅连续，而且有一阶连续导数，于是
\[
 \varphi(x_1)-\varphi(x_2)=\varphi'(\xi)(x_1-x_2),\qquad x_1\le\xi\le x_2.
\]
如果 $\sup|\varphi'(\xi)|=K<+\infty$，则对任给的 $\varepsilon$，令 $\delta(\varepsilon)=\varepsilon/K$，即知当 $|x_1-x_2|<\delta(\varepsilon)$ 时 $|\varphi(x_1)-\varphi(x_2)|<\varepsilon$，而如果这个 $K$ 不只适合某一个 $\varphi(x)$，而且适用于 $M$，则 $M$ 是等度连续的。因此，我们不必假设 $M$ 等度连续，而是换一个空间：
\[
 C^1(I)=\{\varphi(x),\ \varphi\text{ 与 }\varphi'\text{ 在 }[a,b]\text{ 上连续}\}.
\]
在 $C^1(I)$ 中我们可以定义范数为
\[
 \|\varphi\|_{C^1}=\sup\bigl(|\varphi(x)|+|\varphi'(x)|\bigr).
\]
很容易看出 $C^1(I)$ 成了一个巴拿赫空间（即完备的赋范线性空间）。这个空间中的有界集
\[
 M=\{\varphi(x)\in C^1(I),\ \sup(|\varphi(x)|+|\varphi'(x)|)\le K\}
\]
中的元，不但其本身有公共的界 $K$，而且一阶导数也有公共的界 $K$，所以是 $C(I)$（注意，不是 $C^1(I)$ 而是 $C(I)$）中的有公共的界（不妨称为一致有界）的等度连续函数集，因此在 $C(I)$ 中为预紧。$C^1(I)$ 作为一个集合是 $C(I)$ 的子集合，我们就说 $C^1(I)$ 可以嵌入在 $C(I)$ 中。同样上面的 $M$ 也嵌入在 $C(I)$ 中，但是 $M$ 在 $C^1(I)$ 只是 $C^1(I)$ 范数下的有界集，而在 $C(I)$ 中则它——或者说是它的嵌入像——是一个预紧集。我们就说 $M$ 是紧嵌入在 $C(I)$ 中。所以，为了应用 A--A 定理，我们先不妨把 $M$ 放在 $C^1(I)$ 中，并证明它是 $C^1(I)$ 中的有界集。就是说先不妨去证明 $\varphi(x)\in M$ 都有一阶连续导数，然后去估计 $\varphi(x)$ 及其导数 $\varphi'(x)$ 并证明 $\sup(|\varphi(x)|+|\varphi'(x)|)\le K$，$K$ 适用于整个 $M$ 而不只是 $M$ 中某一具体函数。下一步就是证明 $M$ 可以紧嵌入于 $C(I)$ 中，这样从 $M$ 中的一个序列 $\{\varphi_k(x)\}$ 找出一个收敛子序列。因此 A--A 定理就是一个紧嵌入定理。我们只讲了一个最简单的例子，但是读者们应注意到，这是数学和数学物理中极重要的方法。例如在第三章变分学一节中我们讲到某个泛函（例如能量泛函）的极小化序列，它可能是某个索波列夫空间 $H^k(\Omega)$ 中某个集合 $M$ 的序列。我们不知道它是否会收敛，但不妨证明它在另一个空间 $H^l(\Omega)$（$l<k$，可能还要附加一些别的条件）中是紧的，即 $H^k(\Omega)$ 的某个集合 $M$ 紧嵌入于 $H^l(\Omega)$ 中，因此这个极小化序列至少在 $H^l(\Omega)$ 中有收敛的子序列。当然与上面的 $C(I)$ 和 $C^1(I)$ 不同，这里的讨论要以 $L^2$ 理论为基础，但其基本思想是一致的。它们也有共同的问题：上面我们虽然涉及 $C^1(I)$，我们得到的子序列都只是在 $C(I)$ 中收敛，其极限函数只能在 $C(I)$ 中而不在 $C^1(I)$ 中。同样，在变分问题中我们得到的极小化序列也只在 $H^l(\Omega)$ 中，而不在 $H^k(\Omega)$ 中。这里产生的新问题当然又需要其它方法来解决。但不论如何，我们看到了紧性问题是整个数学中的一个重大问题，而决不只是怎样把微积分的基本概念弄得更清楚的问题。

关于紧性的讨论至此为止。
